% mainfile: Refinement.tex
%%% =============================================== Preamble BEGINN =====================================================================
%% Lade:
% Dokumentenklasse
\documentclass[
	draft,	% Zeigt nur Rahmen f"ur Bilder an
	english,%german,
	12pt,		% Schriftgr"o"se 11 Arial oder 12 Times New Roman
	%fleqn,	% fleqn für Matheumgebungen linksb"undig
		% Layout:
	open=right,    % vor dem chapter eine leere linke Seite
	twoside,	% zweiseitiges
	a4paper,	% Papierformat
	%normalheadings,
	%DIV=calc, % berechnet mit \typearea[current]{calc} den richtigen Seitendivisor für die Rasterung der Satzspiegels
	%DIV11,
	%BCOR5mm % Bindekorrektur, was weiß bleibt beim druck und nicht für die Berechnung des Satzspiegels benützt wird
	%left=3cm, % max 4
	%right=3cm, % max 4
	%top=2cm, % max 3
	%bottom=2cm % max 3,
	enabledeprecatedfontcommands
] {scrreprt}		% Art des Dokumentes

% Pakete
% for better fitting on a4 paper
%\usepackage{a4} % kein guter Stil
% Umlaute erm"oglichen
\usepackage[utf8]{inputenc}
% Erm"oglicht beim markieren/kopieren des Textes auch die Umlaute
% au"serdem funktioniert damit auch das Trennen von W"ortern mit Umlauten
% Dadurch aber sehr h"assliche Schrift und Schriftbild, aber mit \usepackage{lmodern}
% l"asst es sich wieder heile machen
\usepackage[T1]{fontenc}
% bindet neue Schriften ein und l"asst das Dokument wieder vom Schriftbild normal aussehen
% was durch \usepackage[T1]{fontenc} zerst"ort wurde
\usepackage{lmodern}
% Neue Rechtschreibung %english f"ur Bibliographie

% Index
\usepackage{index}
% Berechnungen mit Befehlen
%\usepackage{calc}
% deutsches Literaturverzeichnis
%\usepackage{babelbib}
% sprachabh"angiges Datum
\usepackage[english,ngerman]{isodate}
% Einzelne Seiten landscape anzeigen
%\usepackage{lscape}
% Tabellen rotieren k"onnen
%\usepackage{rotating}
% F"ur etwas besser einstellbare Listen
%\usepackage{enumerate} kann weniger als enumitem
%\usepackage{enumitem}

% letztes Argument legt die Sprache des Dokuments fest
%\usepackage[english,ngerman]{babel}
\usepackage[english]{babel}
% Eigentlich veraltet, aber babel kann anscheinend keine richtige Silbentrennung
\usepackage{ngerman}
% F"ur Seitenlayout und Darstellung (vorallem eigener Seitenstil) und header und footer
\usepackage[automark,headsepline]{scrpage2}

% f"ur den Zeilenabstand
\usepackage{setspace}
% f"ur farbigen Text
\usepackage{color}
\usepackage{xcolor}
% Times Roman (Microsoft's Times New Roman)
%\usepackage{txfonts}
% Helvetica (Microsoft's Arial)
%\usepackage{helvet}
% Matheumgebungen
%\usepackage{amsmath}
\usepackage{amssymb}
%\usepackage{amsfonts}
%\usepackage[mathcal]{euscript}
%\usepackage{mathrsfs}
%\usepackage{MnSymbol} % attention overwrites many math symbols for example the frown is very short an object with ^ appear very low.
%% for xLeftarrow and so on and mathclap for spaces
\usepackage{mathtools}
% for a fat ;
\usepackage{stmaryrd}
% f"ur nicefrac
\usepackage{units}
% der Shuffle Operator
\usepackage{shuffle}
%\usepackage{amsthm}
% f"ur die theoreme
\usepackage[amsmath,amsthm,thmmarks]{ntheorem}

% Bilder
%% zwei Bilder in einer Umgebung
%\usepackage{subfig} putt
%\usepackage{subfigure} putt
%\usepackage{subcaption}
% Graphiken einbinden
%\usepackage{graphicx}

% Tabellen
% Um zeilenweise Tabellen zu formatieren
%\usepackage{hhline}
% Tabellen mit einstellbarer Breite
%\usepackage{tabularx}
% Tabellen "uber mehrere Seiten
%\usepackage{longtable}
% passt longtable so an, dass auch die Eigenschaften von tabularx verwandt werden k"onnen
%\usepackage{ltxtable}
% Entzerrt Tabellenzeilen
%\usepackage{booktabs}
% Tabellen farblich markieren
%\usepackage{colortbl}
% Figuren und tables mit "H" exakt an einen Platz binden
\usepackage{float}


% Glossar
%\usepackage[ngerman]{translator}
\usepackage[english]{translator}
\usepackage[
	nowarn,       %stop glossar warninings
	nonumberlist, %keine Seitenzahlen anzeigen
	acronym,      %ein Abkürzungsverzeichnis erstellen
	toc]          %Einträge im Inhaltsverzeichnis
{glossaries}

%Umgebung für Code-Fragmente
\usepackage{listings} \lstset{	
	morecomment=[s][\color{darkred}]{"}{"},
	escapeinside={\%*}{*)},
	commentstyle=\color{cdc_Green},
	keywordstyle=\color{blue},
	frame=single,
	backgroundcolor=\color{lstColor},
	breaklines=true,
	tabsize=4
}


% f"ur Hyperlinks
\usepackage[
	plainpages=false,
	pdfpagelabels,
	breaklinks=true,
	a4paper,
	bookmarks,
	bookmarksopen=true,
	bookmarksnumbered=true,
	pdfauthor={Muhammad Ekbal Ahmad},
	pdftitle={Transformational semantics of the combination pi-OZ for mobile processes with data},
]{hyperref}

% f"ur for-Schleifen
%\usepackage{forloop}
% f"ur if-then-else
\usepackage{ifthen}
% for margins foot header titlepage
\usepackage[a4paper]{geometry}

% TODO notes
%\usepackage[disable]{todonotes} % notes not showed
\usepackage[draft]{todonotes}   % notes showed

% Lorem ipsum blindtext \lipsum or lipsum[n]
\usepackage{lipsum}

% Block comments
%\usepackage{comment}
\usepackage{verbatim}

% Seitenlinien
%\usepackage{framed}
%\usepackage{mdframed}

%% Tikz
\usepackage{pgf}
\usepackage{tikz}

% Neue Befehle
% fix \bf issue
\DeclareOldFontCommand{\bf}{\normalfont\bfseries}{\mathbf}
\DeclareOldFontCommand{\sf}{\normalfont\bfseries}{\mathbf}

% Punkte auch bei section
%\renewcommand\l@section{\@dottedtocline{1}{1.5em}{2.3em}}

% buffer commands
%\newcommand\todo[1]{{\marginpar{\color{red}TODO: #1}}}
%\newcommand\todoI[1]{{\color{red}TODO: #1}}
\newcommand\bla{\todo[inline]{Erkl"arung/Beschreibung/F"ulltext}}
%\newcommand\lorem{Lorem ipsum dolor sit amet, consetetur sadipscing elitr, sed diam nonumy eirmod tempor invidunt ut labore et dolore magna aliquyam erat, sed diam voluptua. At vero eos et accusam et justo duo dolores et ea rebum. Stet clita kasd gubergren, no sea takimata sanctus est Lorem ipsum dolor sit amet.}

% partition commands
\newcommand{\newpart}[1]{{\color{darkred}\noindent\rule{0.5\textwidth-1cm}{1pt} #1 new \rule{0.5\textwidth-1cm}{1pt}}\newline}
\newenvironment{old}[1]
{
	\ifthenelse{\boolean{show_all}}{
		\begin{color}{darkgreen}
		\par\noindent\rule{\textwidth}{1pt}
		\begin{center}\underline{#1}\end{center}
	}{\comment{}}
}
{
	\ifthenelse{\boolean{show_all}}{
		\par\noindent\rule{\textwidth}{1pt}\newline
		\end{color}
	}{\endcomment{}}
}
\newenvironment{new}{
%   \begin{mdframed}[outerlinewidth=2,leftmargin=10,%
%     rightmargin=-10pt,backgroundcolor=white,hidealllines=true,leftline=true,%
%     innertopmargin=0pt,splittopskip=\topskip,skipbelow=\baselineskip,innerbottommargin=0pt%
%     skipabove=\baselineskip]%
	%\begin{leftbar}
%\def\FrameCommand{\vrule width 1pt \hspace{1cm}}%
%\MakeFramed {\advance\hsize\width \FrameRestore}
	\marginpar{\color{darkred}\rule{0.5cm}{1pt} $\downarrow$new \rule{0.5cm}{1pt}}
} {
	\marginpar{\color{darkred}\rule{0.5cm}{1pt} $\uparrow$new \rule{0.5cm}{1pt}}
%\endMakeFramed
	%\end{leftbar}
	%\end{mdframed}
}
% Notizen die bei druckversion nicht angezeigt werden
\newcommand\note[1]{\ifthenelse{\boolean{print_media}}{}{{\color{darkgreen}<Note: #1>}}}

% Referenzieren
\newcommand{\refDef}[1]{Definition~\ref{#1}}
\newcommand{\refFig}[1]{Figure~\ref{#1}}
\newcommand{\refSec}[1]{Section~\ref{#1}}
\newcommand{\refChap}[1]{Chapter~\ref{#1}}
\newcommand{\refLem}[1]{Lemma~\ref{#1}}
\newcommand{\refTheo}[1]{Theorem~\ref{#1}}
\newcommand{\refConv}[1]{Convention~\ref{#1}}
\newcommand{\refConj}[1]{Conjecture~\ref{#1}}
\newcommand{\refEq}[1]{Equation~\ref{#1}}
\newcommand{\refEx}[1]{Example~\ref{#1}}
\newcommand{\refCor}[1]{Corollary~\ref{#1}}
% Referenzieren auf Textstelle
%\newcommand\fullRef[1]{Abschnitt \ref{#1} \glqq\nameref{#1}\grqq{} auf Seite~\pageref{#1}}
%% Referenz auf Bild
%\newcommand\imgRef[1]{Abbildung \ref{#1} auf Seite~\pageref{#1}}
%% Referenz auf Tabelle
%\newcommand\tabRef[1]{Tabelle \ref{#1} auf Seite~\pageref{#1}}
%% Referenz auf Definition
%\newcommand\defRef[1]{Definition \ref{#1} auf Seite~\pageref{#1}}

%Funktionale Anforderungen
%\newcounter{functCounter}
%\newcounter{subFunctCounter}[functCounter]
%\renewcommand\thesubFunctCounter{\Alph{functCounter}\arabic{subFunctCounter}}

%\newcommand\functRequirementPackage{
%\refstepcounter{functCounter}
%\setcounter{subFunctCounter}{0}
%}
%\newcommand\functRequirement[4][\empty]{%%% \empty: Standardwert des optionalen Parameters
%\refstepcounter{subFunctCounter}
%\begin{description}
%\itemsep-0.7cm
%\item[\thesubFunctCounter:] #2 \label{fa_\thesubFunctCounter}
%\item[Eingabe:] #3
%\item[Ausgabe:] #4
%\ifthenelse{\NOT\equal{#1}{\empty}}%%
%		{\item[Optional:] #1}%%
%{}
%\end{description}
%}

% Absatz
%\newcommand\abs{\par\medskip}

% Fetter Index
\newcommand{\findex}[2][\empty]{%%% \empty: Standardwert des optionalen Parameters
	\ifthenelse{\equal{#1}{\empty}}%%
		{\index{#2}\emph{#2}}%%
		{\index{#1}\emph{#2}}%%
	}

% In enumerate-Umgebung (a) solche item
\newcommand{\alphaEnums}{\renewcommand{\labelenumi}{(\alph{enumi})}}

%% Mathematik
\newcommand{\nin}{\not\in}
% Elementzeichen mit anschlie"sendem Mengenzeichen
\newcommand\isin{\in \mathbb}
% Potenzmenge
\newcommand{\pom}[1]{\mathbb P\left(#1\right)}
% Komplexe Zahlen
\newcommand\K{\mathbb C}
% Reelle Zahlen
\newcommand\R{\mathbb R}
% Rationale Zahlen
\newcommand\Q{\mathbb Q}
% Ganze Zahlen
\newcommand\Z{\mathbb Z}
% Nat"urliche Zahlen
\newcommand\N{\mathbb N}
% Kaligraphische Zeichen
\newcommand\calN{\mathcal N}
\newcommand\calL{\mathcal L}
%% backard-Matrix
%\newcommand\B{\mathbb B}
%% forward-Matrix
%\newcommand\F{\mathbb F}
% Menge
\newcommand{\set}[2][\empty]{%%% \empty: Standardwert des optionalen Parameters
	\ifthenelse{\equal{#1}{\empty}}
		{\left\{#2\right\}}
		{\left\{#2\; \mid \; #1 \right\}}
}
\newcommand{\pair}[2]{(#1,#2)}
  
\newcommand{\setmulti}[2][\empty]{%%% \empty: Standardwert des optionalen Parameters
	\ifthenelse{\equal{#1}{\empty}}
		{\bigl\{#2\bigr\}}
		{\bigl\{#2\; \mid \; #1 \bigr\}}
}
\newcommand{\card}[1]{\left|#1\right|}
% Betrag
%\newcommand{\betrag}[1]{\left\lvert #1 \right\rvert}
% Das zu zeigen - Zeichen ZZ
%\newcommand{\zz}{Z\kern-.3em\raise-0.5ex\hbox{Z}:}
% Kongruenz
%\newcommand{\kong}[3]{#1\equiv#2\; mod \;#3}
%\newcommand{\inkong}[3]{#1\nequiv#2\; mod \;#3}
% Vektoren
%\newcommand{\vect}[1]{
%	\left(\begin{array}{c}
%			#1
%		\end{array}
%	\right)
%}
% Matritzen
%\newcommand{\mat}[2]{
%	\left(\begin{array}{#1}
%			#2
%		\end{array}
%	\right)
%}

%% Petri-Netz
% Pre-/Postset
\newcommand{\preset}[1]{{^\bullet{}#1}}
\newcommand{\postset}[1]{{#1^\bullet{}}}

% scalable arrows
\makeatletter
\def\slashedarrowfill@#1#2#3#4#5{%
  $\m@th\thickmuskip0mu\medmuskip\thickmuskip\thinmuskip\thickmuskip
   \relax#5#1\mkern-7mu%
   \cleaders\hbox{$#5\mkern-2mu#2\mkern-2mu$}\hfill
   \mathclap{#3}\mathclap{#2}%
   \cleaders\hbox{$#5\mkern-2mu#2\mkern-2mu$}\hfill
   \mkern-7mu#4$%
}
% already defined in mathtools
%\makeatother
%\makeatletter
%\renewcommand{\xLeftrightarrow}[2][]{\ext@arrow 0359\Leftrightarrowfill@{#1}{#2}}
%\makeatother
%\makeatletter
%\renewcommand{\xleftrightarrow}[2][]{\ext@arrow 0359\leftrightarrowfill@{#1}{#2}}
%\makeatother
%\makeatletter
%\renewcommand{\xLeftarrow}[2][]{\ext@arrow 0359\Leftarrowfill@{#1}{#2}}
%\makeatother
%\makeatletter
%\renewcommand{\xRightarrow}[2][]{\ext@arrow 0359\Rightarrowfill@{#1}{#2}}
%\makeatother
%\makeatletter
\newcommand*{\simfill@}{\arrowfill@\cdot\sim\succ}
\newcommand{\xsimrightarrow}[2][]{\ext@arrow 0359\simfill@{#1}{#2}}
\makeatother
%% negations
\makeatletter
\newcommand*{\nrightarrowfill@}{\slashedarrowfill@\relbar\relbar{\raisebox{.12em}{\tiny/}}\rightarrow}
\newcommand{\nxrightarrow}[2][]{\ext@arrow 0359\nrightarrowfill@{#1}{#2}}
\makeatother

\newcommand{\functText}[1]{\mathtt{#1}}

\newcommand{\ebnf}{\;\bigm|\;}
\newcommand{\gdw}{\mathrm{iff}}
\newcommand{\falls}{\mathrm{if}}
\newcommand{\eq}[1]{\stackrel{#1}{=}}
\newcommand{\pref}[1]{\functText{pref}(#1)}
\newcommand{\conj}[1]{\overline{#1}}
%\newcommand{\true}{\functText{true}}
%\newcommand{\false}{\functText{false}}

%% pi-calculus
\newcommand{\names}{\calN}
\newcommand{\relNames}{\mathfrak{N}}
\newcommand{\conames}{\overline{\calN}}
\newcommand{\labels}{\calL}
\newcommand{\actions}{\functText{Act}}
\newcommand{\outA}{\functText{Out}}
\newcommand{\inA}{\functText{In}}
\newcommand{\boutA}{\functText{Bout}}
\newcommand{\picalc}{$\pi$-calculus}
\newcommand{\syntdef}{::=}
\newcommand{\define}{=_{\mathrm{def}}}
\newcommand{\fnF}{\functText{fn}}
\newcommand{\bnF}{\functText{bn}}
\newcommand{\fn}[1]{\functText{fn}(#1)}
\newcommand{\bn}[1]{\functText{bn}(#1)}
\newcommand{\nF}{\functText{n}}
\newcommand{\n}[1]{\functText{n}(#1)}
\newcommand{\subF}{\functText{sub}}
\newcommand{\objF}{\functText{obj}}
\newcommand{\bindF}{\functText{bind}}
\newcommand{\sub}[1]{\functText{sub}(#1)}
\newcommand{\obj}[1]{\functText{obj}(#1)}
%\newcommand{\objbn}[1]{\functText{obj}_\functText{bn}(#1)}
\newcommand{\bind}[2]{\functText{bind}(#1,#2)}
\newcommand{\bnsubstF}{\functText{bn}_{\functText{subst}}}
\newcommand{\bnsubst}[3]{\functText{bn}_{\functText{subst}}(\subs{#1}{#2},#3)}
\newcommand{\struc}{\equiv{}}
\newcommand{\struct}[2]{#1\struc{}#2}
% sequences
\newcommand{\seqset}[1]{\functText{seq}(#1)}
\newcommand{\eseq}{\langle\rangle}
%\newcommand{\seq}[1]{\langle#1\rangle}
\newcommand{\seqconc}[2]{{#1}^{\,\frown{}\,}{#2}}
\newcommand{\seqcom}[2]{#1\leftrightharpoons{}#2}
\newcommand{\subsetsim}{\mathrel{\substack{\textstyle\subset\\[-0.5ex]\textstyle\sim}}}
\newcommand{\simcirc}{\mathrel{\substack{\textstyle\sim\\[-0.4ex]\textstyle\circ}}}
\newcommand{\ec}[1]{\left[#1\right]}%_{\alpha}}
\newcommand{\bigstep}[1]{\xRightarrow{#1}}
\newcommand{\len}[1]{\#(#1)}
%\newcommand{\parl}[1]{\overrightarrow{#1}}
\newcommand{\parl}[1]{\vec{#1}}
\newcommand{\substF}{\sigma}
\newcommand{\subst}[1]{\left(#1\right)\sigma}
%\newcommand{\subs}[2]{\left\{\nicefrac{#1}{#2}\right\}}
\newcommand{\substitue}[2]{\left\{\nicefrac{#1}{#2}\right\}}

% prefix
\newcommand{\singleout}[1]{\bar{#1}}
\newcommand{\out}[2]{\overline{#1}\langle#2\rangle}
\newcommand{\outa}[2]{\overline{#1}\langle#2\rangle}
\newcommand{\bout}[2]{\overline{#1}(#2)}
\newcommand{\bouta}[2]{\overline{#1}(#2)}
\newcommand{\inp}[2]{#1(#2)}
\newcommand{\inpa}[2]{#1\,#2}
% processes
\newcommand{\procs}{\mathcal P^{\pi}}
\newcommand{\sums}{\procs_{M}}
\newcommand{\procsesf}{\procs_{\functText{esf}}}
\newcommand{\procsresf}{\procs_{\functText{resf}}}
\newcommand{\procsrecf}{\procs_{\functText{recf}}}
\newcommand{\procchoice}[2]{#1 + #2}
\newcommand{\procsum}{\sum_{i\in I}{\pi_i.P_i}}
\newcommand{\procpar}[2]{#1 \mid #2}
%\newcommand{\procres}[3][\empty]{%%% \empty: Standardwert des optionalen Parameters
%	\ifthenelse{\equal{#1}{\empty}}{
%		\ifthenelse{\equal{#2}{\empty}}{
%			\textnormal{\underline{\texttt{new}}}
%		}{
%			\textnormal{\underline{\texttt{new}}}\,#2\;#3
%		}
%	}{
%		\textnormal{\underline{\texttt{new}}}\,#2\left(#3\right)
%	}
%}
\newcommand{\procres}[3][\empty]{%%% \empty: Standardwert des optionalen Parameters
	\ifthenelse{\equal{#1}{\empty}}{
		\ifthenelse{\equal{#2}{\empty}}{
			\underline{\functText{new}}
		}{
			\underline{\functText{new}}\,#2\;#3
		}
	}{
		\underline{\functText{new}}\,#2\left(#3\right)
	}
}
\newcommand{\proccall}[2]{#1\langle#2\rangle}
\newcommand{\procdef}[2]{#1(#2)}
\newcommand{\proczero}{\mathbf{0}}
% substitution
\newcommand{\transp}[2]{\left\{#1\leftrightarrow{}#2\right\}}% substitution
\newcommand{\transpT}[2]{\theta_{#1}(#2)}
\newcommand{\supp}[1]{\functText{supp}(#1)}
\newcommand{\cosupp}[1]{\functText{cosupp}(#1)}
% alpha conversion
\newcommand{\alphaeq}{=_\alpha}
% transitions
%% sangiorgi
\newcommand{\transs}[1]{\xrightarrow{#1}}
\newcommand{\intrans}[2]{\transs{\inpa{#1}{#2}}}
\newcommand{\outtrans}[2]{\transs{\out{#1}{#2}}}
\newcommand{\bouttrans}[2]{\transs{\bout{#1}{#2}}}
\newcommand{\tautrans}{\transs{\tau}}
%% mine
\newcommand{\trans}[1]{\xsimrightarrow{#1}}
%\newcommand{\trans}[1]{\rightsquigarrow}
\newcommand{\transin}[2]{\trans{\inp{#1}{#2}}}
\newcommand{\transout}[2]{\trans{\out{#1}{#2}}}
\newcommand{\transbout}[2]{\trans{\bout{#1}{#2}}}
\newcommand{\transtau}{\trans{\tau}}
% rules name, pr"amisse1, pr"amisse2, konklusion, optional Anwendungsbedingung
\newcommand{\kalRule}[5][\empty]{%%% \empty: Standardwert des optionalen Parameters
	\underline{\scriptstyle#2}:\;
	\ifthenelse{\equal{#3}{\empty}}{
		\ifthenelse{\equal{#4}{\empty}}{
			\dfrac{}{#5}
		}{
			\dfrac{#4}{#5}
		}
	}{
		\dfrac{#3 \quad #4}{#5}
	}
	\ifthenelse{\NOT\equal{#1}{\empty}}{\;\;#1}{}
}
% names of rules
\newcommand{\rulename}[1]{$#1$}
\newcommand{\etau}{\rulename{E-TAU}}
\newcommand{\ecall}{\rulename{E-CALL}}
\newcommand{\eout}{\rulename{E-OUT}}
\newcommand{\ein}{\rulename{E-IN}}
\newcommand{\esuml}{\rulename{E-SUM_L}}
\newcommand{\esumr}{\rulename{E-SUM_R}}
\newcommand{\eres}{\rulename{E-RES}}
\newcommand{\eparl}{\rulename{E-PAR_L}}
\newcommand{\eparr}{\rulename{E-PAR_R}}
\newcommand{\eopen}{\rulename{E-OPEN}}
\newcommand{\ecoml}{\rulename{E-COM_L}}
\newcommand{\ecomr}{\rulename{E-COM_R}}
\newcommand{\eclosel}{\rulename{E-CLOSE_L}}
\newcommand{\ecloser}{\rulename{E-CLOSE_R}}
% strong / weak Bisimulation
\newcommand{\sbisim}{\sim}
\newcommand{\wbisim}{\approx}
\newcommand{\simu}{\mathcal{S}}
\newcommand{\weaksimuset}[2]{\simu^{#1,#2}}
\newcommand{\simuset}{\weaksimuset{P}{Q}}
% denotational semantics
\newcommand{\tr}{\functText{Traces}}
\newcommand{\trI}[1]{\mathcal{T_I}(\ec{#1})}
\newcommand{\tracesI}[3]{\mathcal T_{#1}(\ec{\procpar{#2}{#3}})}
\newcommand{\traces}[1][\empty]{
		\ifthenelse{\equal{#1}{\empty}}
		{\mathcal T}
		{\mathcal T(\ec{#1})}
}

\newcommand{\tracesR}[1]{\mathcal{T}_\relNames(\ec{#1})}

\newcommand{\res}[3]{\functText{res}(#1,#2,\traces[#3])}
\newcommand{\failures}[1][\empty]{
		\ifthenelse{\equal{#1}{\empty}}
		{\mathcal F}
		{\mathcal F(#1)}
}
\newcommand{\fd}[1][\empty]{
		\ifthenelse{\equal{#1}{\empty}}
		{\mathcal{FD}}
		{\mathcal{FD}(#1)}
}
% refinement
\newcommand{\refi}[1][\empty]{
		\ifthenelse{\equal{#1}{\empty}}
		{\,\sqsubseteq_{\mathcal T}\,}
		{\,\sqsubseteq_{\mathcal{#1}}\,}
}
\newcommand{\rrefi}[1][\empty]{
		\ifthenelse{\equal{#1}{\empty}}
		{\,\sqsupseteq_{\mathcal T}\,}
		{\,\sqsupseteq_{\mathcal{#1}}\,}
}

% Header/Footer
% style
%\pagestyle{scrheadings}
\pagestyle{empty} % first empty, then until introduction scrheadings
% Konfiguration der Kopf- und Fu"szeilen: siehe scrguide.pdf Kapitel 5 (ab Seite 214)
% pagenumber first roman until introduction arabic
\pagenumbering{Roman} 

% Settings
% Variable je nach dem wie sie gesetzt ist, wird die Druckversion beziehungsweise die digitale Version erstellt
% Bedeutet Links farbig oder schwarz
\newboolean{print_media}
\setboolean{print_media}{false}

% Ob alle Teile angezeigt werden sollen
\newboolean{show_all}
\setboolean{show_all}{false}
%\includecomment{myComment}

%\KOMAoption{BCOR}{7mm}
%\KOMAoption{parskip}{yes}
\KOMAoption{toc}{bib,index,listof}

% Abstand itemize "andern
%\setitemize{itemsep=-0.5cm}

% Seitenränder
\geometry{
	includehead,
	includefoot,
	inner=3cm, % min 3 max 4
	outer=3cm, % min 3 max 4
	top=2cm, % min 2 max 3
	bottom=2cm % min 2 max 3
}
%\geometry{showframe}

% Farben
\definecolor{gruen}{rgb}{.10, .4, .50}
\definecolor{blau}{rgb}{.2, .2, 0.6}
\definecolor{grau}{gray}{.75}
%\definecolor{white}{gray}{.75}
\definecolor{lstColor}{rgb}{0.901,0.901,0.901} %% Back (Almost White)
\definecolor{darkred}{rgb}{0.601,0.001,0.001} %% dunkles Rot
\definecolor{darkgreen}{rgb}{.10, .4, .20}% dunkelgruen
\definecolor{cdc_Blue}{rgb}{0.0,0.355,0.652} %% Heading Blue (Oldenburg CD Ultramarine)
\definecolor{cdc_BlueM}{rgb}{0.25,0.473,0.722} %% Heading Blue Medium (Oldenburg CD Ultramarine 60%)
\definecolor{cdc_BlueL}{rgb}{0.5,0.589,0.793} %% Heading Blue Light (Oldenburg CD Ultramarine 20%)
\definecolor{cdc_Green}{rgb}{0.390,0.695,0.285} %% Heading Green (Oldenburg CD Chartreuse)
\definecolor{cdc_GreenM}{rgb}{0.559,0.758,0.438} %% Heading Green Medium (Oldenburg CD Chartreuse 60%)
\definecolor{cdc_GreenL}{rgb}{0.688,0.828,0.605} %% Heading Green Light (Oldenburg CD Chartreuse 20%)
\definecolor{color_back}{rgb}{0.941,0.941,0.941} %% Back (Almost White)

% Zeilenabstand
\onehalfspacing
% Index
%\newindex{default}{idx}{ind}{Index}
%\newindex{name}{adx}{and}{Namensverzeichnis}
%\newindex{defi}{ddx}{dnd}{Definitionsverzeichnis}
%\newindex{satz}{sdx}{snd}{Satzverzeichnis}
%\newindex{lem}{ldx}{lnd}{Lemmataverzeichnis}
%\newindex{kor}{kdx}{knd}{Korollarverzeichnis}
%\newindex{bemerkung}{bdx}{bnd}{Bemerkungsverzeichnis}
%\newindex{beispiel}{bsdx}{bsnd}{Beispielverzeichnis}

%Einstellungen Code-Fragmente
%\lstset{
%  xleftmargin=13pt,
%  xrightmargin = 5pt,
%  basicstyle=\small\ttfamily,
%  columns=fullflexible,
%  showstringspaces=false,
%  commentstyle=\color{gray}\upshape,
%  literate={"a}{\"a}1 {"o}{\"o}1 {"u}{\"u}1 {"s}{\ss}1
%  {"A}{\"A}1 {"O}{\"O}1 {"U}{\"U}1
%}

%\lstloadlanguages{Java }

%% Listingstyles
\lstdefinelanguage{pseudocode}
{
	alsolanguage=Java,
	morekeywords=[1]{INPUT, OUTPUT, then, od, in, fi, foreach, begin, end, endfor, endforeach, endwhile},
	morecomment=[s][\color{darkred}]{"}{"},
	escapeinside={\%*}{*)},
	commentstyle=\color{cdc_Green},
	keywordstyle=\color{blue},
	frame=single,
	backgroundcolor=\color{white},
	xleftmargin=5pt,
	xrightmargin = 5pt,
	breaklines=true,
}
%\lstdefinelanguage{apt}
%{
%	alsolanguage=Java,
%	morekeywords=[2]{Place, Transition, Arc, Flow, PetriNet, TransitionSystem, Marking, Node, Token, ArcKey, State, Edge, EdgeKey, IEdge, IGraph, INode, PetriNetOrTransitionSystem},
%	keywordstyle=[2]\color{cdc_BlueM},
%	morecomment=[s][\color{darkred}]{"}{"},
%	commentstyle=\color{cdc_Green},
%	keywordstyle=\color{blue},
%	escapeinside={\%*}{*)},
%	frame=single,
%	%backgroundcolor=\color{lstColor},
%	breaklines=true,
%	xleftmargin=5pt,
%	xrightmargin = 5pt,
%	literate= {Ö}{{\"O}}1 {Ä}{{\"A}}1 {Ü}{{\"U}}1 {ß}{{\ss}}2 {ü}{{\"u}}1 {ä}{{\"a}}1 {ö}{{\"o}}1,
%	extendedchars=true
%}

%\lstdefinelanguage{apt-parser}
%{
%	alsolanguage=Java,
%	morekeywords=[2]{IPNParserOutput, AptPNFormatParser, AptPNFormatLexer,APTPNParserOutput,ParserContext, APTParserContext, PetriNet, ANTLRParser, TransitionSystem, SynetLTSParser, SynetPNParser, APTLTSParser, APTPNParser, APTParser,APTRenderer},
%	keywordstyle=[2]\color{cdc_BlueM},
%	morekeywords=[3]{CommonTokenStream, RecognitionException},
%	keywordstyle=[3]\color{cdc_BlueL},
%	morecomment=[s][\color{darkred}]{"}{"},
%	commentstyle=\color{cdc_Green},
%	keywordstyle=\color{blue},
%	escapeinside={\%*}{*)},
%	frame=single,
%	%backgroundcolor=\color{lstColor},
%	breaklines=true,
%	xleftmargin=5pt,
%	xrightmargin = 5pt,
%	literate= {Ö}{{\"O}}1 {Ä}{{\"A}}1 {Ü}{{\"U}}1 {ß}{{\ss}}2 {ü}{{\"u}}1 {ä}{{\"a}}1 {ö}{{\"o}}1,
%	extendedchars=true
%}

%\lstdefinelanguage{apt-format}
%{
%	morekeywords=[1]{.name, .type, .description, .transitions, .places, .flows, .initial_marking, .final_markings,
%			.states, .labels, .arcs},
%	alsoletter={.},
%	keywordstyle=[1]\color{cdc_BlueM},
%	morecomment=[s][\color{darkred}]{"}{"},
%	morecomment=[s][\color{cdc_Green}]{/*}{*/},
%	morecomment=[l][\color{cdc_Green}]{//},
%	escapeinside={\%*}{*)},
%	frame=single,
%	%backgroundcolor=\color{lstColor},
%	breaklines=true,
%	xleftmargin=5pt,
%	xrightmargin = 5pt,
%	numbers=none,
%}

%\lstdefinelanguage{ebnf}
%{
%	morecomment=[s][\color{darkred}]{'}{'},
%	escapeinside={\%*}{*)},
%	morecomment=[s][\color{cdc_Green}]{(*}{*)},
%	morecomment=[s][\color{cdc_BlueM}]{?}{?},
%	frame=single,
%	%backgroundcolor=\color{lstColor},
%	breaklines=true,
%	xleftmargin=5pt,
%	xrightmargin = 5pt,
%	showspaces=false,
%	numbers=none  
%}

%\lstdefinelanguage{my_xml}
%{
%	morestring=[b]",
%	morestring=[s]{>}{<},
%	morecomment=[s][\color{darkred}]{"}{"},
%	morecomment=[s]{<?}{?>},
%	stringstyle=\color{black},
%	identifierstyle=\color{blue},
%	keywordstyle=\color{cdc_BlueM},
%        morekeywords={classpathref,classname,fork, failonerror, value, path}
%	xleftmargin=5pt,
%	xrightmargin = 5pt,
%	numbers=none,
%}

%\lstdefinelanguage{synet}
%{
%	morecomment=[s][\color{darkred}]{'}{'},
%	escapeinside={\%*}{*)},
%	frame=single,
%	%backgroundcolor=\color{lstColor},
%	breaklines=true,
%	xleftmargin=5pt,
%	xrightmargin = 5pt,
%	numbers=none
%}

% Mathematik Umgebungen
\makeatletter
\newtheoremstyle{plainWithSeparator}
	{\item[\hskip\labelsep \theorem@headerfont ##1\ ##2\theorem@separator]}%
	{\item[\hskip\labelsep \theorem@headerfont ##1\ ##2\ (##3)\theorem@separator]}
\makeatother
\makeatletter
\newtheoremstyle{nonumberPlainWithSeparator}
	{\item[\hskip\labelsep \theorem@headerfont ##1\theorem@separator]}%
	{\item[\hskip\labelsep \theorem@headerfont ##1\ (##3)\theorem@separator]}
\makeatother
\theoremstyle{plainWithSeparator}
\theoremheaderfont{\normalfont\bfseries}
\theorembodyfont{\itshape}
\theoremseparator{}
%\theoremindent{0.5cm}%0.8cm

% theorems
\theoremsymbol{\ensuremath{\square}}
\newtheorem{theorem}{Theorem}[section]

% Lemmata
\theoremsymbol{\ensuremath{\diamondsuit}}
\newtheorem{lemma}{Lemma}[section]

% conventions
\theoremsymbol{\ensuremath{\circ}}
\newtheorem{conv}{Convention}[section]

% conjectures
\theoremsymbol{\ensuremath{\circ}}
\newtheorem{conject}{Conjecture}[section]

%% corollaries
\theoremsymbol{\ensuremath{\square}}
\newtheorem{cor}{Corollary}[section]

%% corollaries
\theoremsymbol{\ensuremath{\square}}
\newtheorem{rem}{Remark}[section]

%%Propositions
%\theoremsymbol{\ensuremath{\square}}
%\newtheorem{proposition}{Proposition}[section]

% examples
\theorembodyfont{\upshape}
\theoremsymbol{\ensuremath{\ast}}
\newtheorem{example}{Example}[section]

% Bemerkung
%\theoremsymbol{}
%\newtheorem{bemerkung}{Bemerkung}[section]

% Definitionen
\theoremsymbol{\ensuremath{\triangle}}
\newtheorem{definition}{Definition}[section]

% Beweise
%\theoremstyle{plain}
\theoremstyle{nonumberPlainWithSeparator}
\theoremseparator{:}
\theoremindent0cm
\theoremsymbol{\rule{1.5ex}{1.5ex}}
\newtheorem{prf}{Proof}

% Hyperlinks
\ifthenelse{\boolean{print_media}}{
	\hypersetup{
		colorlinks=false
		,unicode=true
		,pdfborder={0 0 0}
		,breaklinks=true
		%,linktocpage % falls dort Umbrueche noetig sind
	}
}{
	\hypersetup{colorlinks
		,linkcolor=blau
		,unicode=true
		,urlcolor=blau
		,citecolor=gruen
		,filecolor=gruen
		,breaklinks=true
		%,linktocpage % falls dort Umbrueche noetig sind
	}
}

% Bibliographie
% urldate-format sprachabh"angig machen
%\setbibliographyfont{urldate}{\printdate}
% Definition des Aussehens der Literaturlisten
\bibliographystyle{alpha}

% Erlaubt gr"o"sere Abst"ande, damit nicht "uber den Rand geschrieben wird
%\emergencystretch=1ex

% Erlaubt den Umbruch von align-Umgebungen (Seitenumbruch)
\allowdisplaybreaks

% TIKZ
\usetikzlibrary{arrows, automata, positioning, decorations.pathmorphing}
\tikzstyle{every state}=[draw=none, rectangle, text=black, yshift=10mm]
\tikzstyle{sdots}=[yshift=5mm]
\tikzstyle{ldiag}=[left=2mm,pos=0.4]
\tikzstyle{rdiag}=[right=2mm,pos=0.4]
\tikzstyle{dopac}=[dashed, opacity=0.35]
\tikzstyle{every picture}=[
				->,
				>=stealth',
				%shorten >=1pt,
				%auto, % position of labels
				node distance=3cm,
                    		semithick, % thickness of arrows
				initial text={},
				initial where=above]

%\typearea[current]{calc} % berechnet den Divisor für das Raster zur Berechnung des Satzspiegels

%%%  =============================================== Preamble END ======================================================================

\textbf{Aufbereitungsdienst:}  Ein Aufbereitungsdienst verändert Umweltdaten, die der IoT-Plattform bereitgestellt werden.\\
\\
\textbf{aufbereitete Daten:} Aufbereitete Daten sind Daten, die das Ergebnis eines Berechnungsverfahren eines Aufbereitungsdienstes sind. \\
\\
\textbf{Akzeptanzkriterium:} Ein Akzeptanzkriterium ist eine Anforderung an eine User-Story, die zum Zeitpunkt der Fertigstellung der User-Story erfüllt sein muss. \\
\\
\textbf{BME280:} Der BME280 ist der Sensor für Luftdruck, Temperatur und ralative Luftfeuchte. \\
\\
\textbf{Datenqualität:}  Unter Datenqualität werden Daten verstanden, die vom Sensor gemessen wurden und nicht durch Einflussfaktoren verfälscht sind.\\
\\
\textbf{Definition of Done:} Eine Definition of Done ist eine Art Checkliste. Erst wenn alle Kriterien der Definition of Done erfüllt sind, darf die Story (Anforderung) als "erledigt" beziehungsweise fertig implementiert angesehen werden. Die Definition of Done wird in jedem Story-Ticket gepflegt. \\
\\
\textbf{Definition of Ready:} Eine Definition of Ready ist auf User Stories bezogen und gibt an, ab wenn eine User Story so gepflegt ist, dass sie "ready" ist, um in einem Sprint bearbeitet zu werden.  \\
\\
\textbf{Dienst:}  Ein Dienst stellt der IoT-Plattform Daten bereit oder fragt Daten ab.\\
\\
\textbf{Dynamisches Routing:}
Dynamisches Routing ist ein Verfahren, in welchem aufgrund von veränderten Umweltdaten die Route während der Navigation neu berechnet wird.\\
\\
\textbf{Dynamische Navigation:}
Dynamische Navigation ist die Anpassung der Leitung des Nutzers entlang der Route. Diese Anpassung erfolgt aufgrund der Nichteinhaltung der Route oder aufgrund der Auswahl einer alternativen Route, die aufgrund des dynamischen Routings zur Verfügung steht.\\
\\
\textbf{Epic:} Ein Epic ist die Abstraktionsebene unter den Super Epics. Sie zerlegt das Super Epic in mehrere Epics, die das jeweilige Super Epic spezifizieren. \\
\\
\textbf{Externer Dienst:} Ein externer Dienst ist ein Dienst, der außerhalb der IoT-Plattform liegt und Daten von dieser abfragt. Das kann z.B. der Routing Dienst sein.\\
\\
\textbf{Event - Benachrichtigung:} Eine Event - Benachrichtigung ist eine Mitteiöung, die dem Nutzer angezeigt wird, wenn z.B. ein Verbindungsverlust auftritt oder der Standort eines Sensorknotens verändert wurde. \\
\\
\textbf{Feinstaub:} Feinstaub ist Staub, der aus kleinen Partikeln besteht, wie z.B. gefährliche Stoffe wie Schwermetall aber auch Staubpartike selbst. Dieser kann verschiedene Größen haben; PM 2.5 und PM 10. \\
\\
\textbf{Firmware:} Unter Firmware wird die Software verstanden, die auf dem Sensorknoten läuft.\\
\\
\textbf{GNSS:} Ein GNSS ist ein Global Navigation Satellite-System. Dazu gehören zum Beispiel "GPS" als amerikanisches Satellitensystem, oder "Galileo" als europäisches System.\\
\\
\textbf{Highcharts:}
Bei Highcharts handelt es sich um ...\\
\\
\textbf{historische Daten:} Historische Daten sind Daten aus früheren Messungen der Sensorknoten.  \\
\\
\textbf{IoT:}
IoT ist die Abkürzung für \dq Internet of Things\dq  und ist ein Sammelbegriff für Technologien einer globalen Infrastruktur der Informationsgesellschaften, die es ermöglicht, physische und virtuelle Gegenstände miteinander zu vernetzen und sie durch Informations-und Kommunikationstechniken zusammenarbeiten zu lassen.\\
\\
\textbf{IoT-Plattform:} Die IoT-Plattform nimmt Daten von Sensoren an, speichert diese ab und stellt sie externen Diensten bereit. Außerdem authentifizieren sich Sensoren externe Dienste und Aufbereitungsdienste mit der IoT-Plattform. \\
\\
\textbf{JIRA:} JIRA ist ein Ticketing-System im Projektmanagement, mithilfe dessen für unterschiedliche Anliegen Tickets erstellt und zugewiesen werden können.  \\
\\
\textbf{JIRA-Label:} Ein Label beschreibt in JIRA ein Stichwort, welches einheitlich für Tickets gepflegt werden kann.  \\
\\
\textbf{JWT:} Ein JSON-Web-Token (kurz JWT) ist ein auf JSON basiertes Zugriffstoken.  \\
\\
\textbf{korrekte Daten:} Korrekte Daten sind die tatsächlich vom Sensor gemessenen Daten.\\
\\
\textbf{Metadaten oder Metainformationen:} Metadaten oder Metainformationen beschreiben die eigentlichen Daten, indem sie Informationen über Merkmale der Daten enthalten, wie z.B. Messverfahren, ID oder den Standort. \\
\\
\textbf{MongoDB:} MongoDB ist eine dokumentenorientierte NoSQL-Datenbank, mit Hilfe derer man unter anderem JSON-ähnliche Dokumente verwalten kann.  \\
\\
\textbf{Navigation:}
Navigation ist die Leitung eines Nutzers entlang einer berechneten Route.\\
\\
\textbf{Navigationsanwendung:} Die Navigationsanwendung ist die Endanwendung, mithilfe derer sich der Nutzer eine Route generieren lassen kann und auf ebendieser Route navigiert wird.  \\
\\
\textbf{Navigationsnutzer:} Der Navigationsnutzer ist der Nutzer der Navigationsanwendung. Er kann sich anhand festgelegter Parameter eine Route generieren lassen und hat die Möglichkeit über die ausgewählte Route navigiert zu werden.  \\
\\
\textbf{nicht - funktionale Anforderungen:} Nicht - funktionale Anforderungen sind Anforderungen an die Qualität der Funktionalität, wie z.B.Performance. \\
\\
\textbf{NodeMCU:} Das NodeMCU-Board ist ein Kommunikations - und Verarbeitungssystem für den Sensorknoten. \\
\\
\textbf{PM 2,5:} Der Feinstaub besteht aus 50 Prozent der Teilchen mit einem Durchmesser von 2.5  $\mu$m.\\
\\
\textbf{PM10:} Der Feinstaub besteht aus 50 Prozent der Teilchen mit einem Durchmesser von 10 $\mu$m. \\
\\
\textbf{Prognose:} Eine Prognose ist eine Vorhersage über die Feinbelastung oder über andere Umweltdaten. \\
\\
\textbf{Rohdaten:} Rohdaten sind die Daten, die nicht durch einen Dienst aufbereitet oder verändert wurden. \\
\\
\textbf{Routing:}
Routing ist ein Verfahren, in welchem aufgrund von Umweltdaten eine Strecke von einem eingegebenen Start- und Endpunkt berechnet wird.\\
\\
\textbf{SDS011:} Der SDS011 ist der Feinstaubsensor.\\
\\
\textbf{Stakeholder:} Ein Stakeholder ist eine Person oder Gruppe, die ein berechtigtes Interesse am Verlauf oder Ergebnis eines Prozesses oder Projektes hat. \\
\\
\textbf{Sensor:} Ein Sensor ist ein technisches Bauteil, das bestimmte physikalische (oder chemische Eigenschaften) seiner Umgebung qualitativ oder als Messgröße quantitativ erfassen kann.\\
\\
\textbf{Sensorknoten:} Ein Sensorknoten ist eine Komposition bestehend aus einem Verarbeitungssystem und mehreren Messsystemen. Dabei besteht das Verarbeitungssystem aus einem Controller und einer zugehörigen Kommunikationseinheit. Das Messsystem sind Sensoren, die die Umweltdaten messen, wie zum Beispiel Luftdruck, Feinstaub oder Temperatur. \\
\\
\textbf{Super Epic:} Ein Super Epic ist die höchste Abstraktionsebene beim Formulieren der Anforderungen. Sie beschreibt grob einen Themenkomplex und ist in mehrere darunterliegende Epics aufgeteilt. \\
\\
\textbf{Unixformat:}
Das Unixformat ist....\\
\\
\textbf{Umweltdaten:} Umweltdaten sind Werte über den Zustand von Umweltbestandteile, die durch Beobachtungen, Messungen u.a. gewonnen wurden.
\\
\textbf{UserStory:} Eine User Story ist eine Nutzergeschichte, die eine funktionale Anforderung an das System aus Sicht eines bestimmten Stakeholders formuliert. Sie zerlegt das Epic in konkrete funktionale Anforderungen.  \\
\\
\textbf{Task:} Eine Task ist ein Tickettyp in Jira, der üblicherweise User Stories zugeordnet wird und die technische Umsetzung der Story beschreibt. Tasks können aber auch unabhängig von User Stories verwendet werden, um Aufgaben zu verteilen.  \\
\\
\textbf{UI:} Die UI (User Interface) stellt die visuelle Schnittstelle zum Nutzer dar. Hier ist insbesondere die Benutzeroberfläche des Systems gemeint, welche durch die Oberfläche der Navigationsanwendung und der Webanwendung abgebildet wird.\\
\\
\textbf{Umweltdaten:} Umweltdaten sind Werte über den Zustand von Umweltbestandteilen, die durch Beobachtungen, Messungen u.a. gewonnen wurden.
Zu den Umweltbestandteilen gehören z.B. Boden, Luft oder Gewässer.\\
\\
\textbf{virtuelle Sensorknoten:} Ein virtueller Sensorknoten ist ein Sensorknoten, der nicht physisch Daten mittels Sensoren misst sondern Daten bereitstellt, die mittels verschiedener Berechnungsverfahren, wie z.B. der Mittelwert, berechnet wurden. \\
 


%%%  =============================================== Document BEGINN ===================================================================
\begin{document}

	%% ------------------------------------------------- Main ----------------------------------------------------------------------
	%% Titel
	\selectlanguage{ngerman}
	% mainfile: Refinement.tex
\def\myTitle{Transformational semantics of the combination \emph{$\pi$-OZ} for mobile processes with data}
%\def\mySubtitle{asdfasdf}
\def\myDocumentType{Masterarbeit}
\def\myName{Muhammad Ekbal Ahmad}
\def\myAddress{Brunsbrok 4}
\def\myPlz{26127}
\def\myCity{Oldenburg}
\def\myEmail{muhammad.ekbal.ahmad@uni-oldenburg.de}
\def\myBirth{05.10.1985}
\def\myUni{Universit"at Oldenburg}
\def\myMatrNr{4565577}
\def\myFakultaet{II}
\def\myFachbereich{Informatik, Wirtschafts- und Rechtswissenschaften}
\def\myInstitut{Department f"ur Informatik}
\def\myStudiengang{Fach-Master Informatik}
\def\myFirstPruefer{Prof. Dr. Ernst-R"udiger Olderog}
\def\mySecondPruefer{M.Sc. Manuel Gieseking}
\def\myBetreuer{}
\def\myAbteilung{Entwicklung korrekter Systeme}

% Zur Seiteneinstellung
\def\myTopMargin{15mm}
%\def\myBottomMargin{50mm}
\def\myEnlargement{30mm}

% Kopf- und Fu"szeilen l"oschen
\thispagestyle{empty}
% Einstellungen der Seitenr"ander
\addtolength{\topmargin}{-\myTopMargin}
% Textheigth + x
\enlargethispage{\myEnlargement}
\vspace*{-2cm}
\begin{figure}[htbp]
	\centering
	%\subfloat{\includegraphics[width=0.5\textwidth]{./chapters/0_title/logo_first.eps}}
	%\quad \subfloat{\includegraphics[width=0.25\textwidth]{./chapters/0_title/logo_second.jpg}}
	\includegraphics[width=0.5\textwidth]{./images/logoUniOL.pdf}
\end{figure}
\begin{center}
	\bf{Fakult"at \myFakultaet: \myFachbereich \newline \myInstitut}

	%\vspace*{-0.6cm}
	\bf{Abteilung: \myAbteilung}    
	\vspace*{0.5cm}

	{\color{grau}\noindent\rule{\textwidth-5cm}{1pt}}

	\vspace*{4.5cm}

	{\bf \huge \myTitle \\}
	\vspace*{1.5cm}
	%   \vspace*{0.5cm}
	%  {\bf \large \mySubtitle}
	% \vspace*{1cm} \\
	{\Large \myDocumentType\\}
	\vspace*{0.5cm}
	{\large \it\texttt{-- post version --}}
\end{center}
\vfill
\normalsize{
	\begin{tabular}{ll}
	    	Name: & {\myName} \\
		%Stra"se: & \myAddress \\
		%Wohnort: & \myPlz \ \myCity \\ 
		%Geburtstag: & \myBirth \\ \\
	    	%Matrikelnr.: & {\myMatrNr} \\
		E-Mail: & \myEmail \\ \\
	    	Studiengang: & \myStudiengang\\
		%Betreuer: & \myBetreuer \\
		Erstgutachter: & \myFirstPruefer \\
		Zweitgutachter: & \mySecondPruefer \\
		Datum: & \today \\
	\end{tabular}
}
% Seitenlayout zur"ucksetzen
\addtolength{\topmargin}{\myTopMargin}

	\selectlanguage{english}
	%\cleardoublepage
	%% mainfile: Refinement.tex
\begin{abstract}
\todo{asdf}
\end{abstract}


	%\cleardoublepage
	%\newpage
	%% Inhaltsverzeichnis
	\singlespacing
	\tableofcontents
	\onehalfspacing
	% Abbildungsverzeichnis
	\listoffigures
	% Tabellenverzeichnis
	%\listoftables	
	%Listingsverzeichnis
	%\renewcommand{\lstlistlistingname}{Listingsverzeichnis}
	%\lstlistoflistings
	%% Main
	% mainfile: ../Refinement.tex
\chapter{Introduction}
\pagestyle{scrheadings}	
\setcounter{page}{0}
\pagenumbering{arabic} 
\label{sec_introduction}
In many cases of modern computing it is of interest to describe and model concurrency. Computers no longer just solve a problem by subsequently working off the single tasks of their own, but they decompose and concurrently calculate the problem even together in a network. The increase in the number of CPU cores and more heavily of GPU cores within one single computer convincingly demonstrates how fundamental concurrency is for modern computing. Moreover, the rapidly increasing spread of the Internet is one of the most common examples which shows the importance of networks.

For concurrently calculating the decomposed tasks, these tasks have to be distributed to the participants. Moreover, the members have to communicate with each other, for example, for sending interrupts and sharing results. Generally, one can say, the more use of concurrency the more communication has to be done. Hence, concurrently communicating systems are rather the norm than an exception. Thus, investigating concurrency and communication are a continuing concern within modern computing.

%One well-accepted way to specify such systems are \findex[process algebra]{process algebras} like Tony Hoare's \gls{CSP} \cite{hoareOriginalCSP} and Robin Milner's \gls{CCS} \cite{milnerCCS}.
For the formal specification of such systems \findex[process!algebra]{process algebras} have been widely accepted, especially Tony Hoare's \gls{CSP} \cite{hoareOriginalCSP} and Robin Milner's \gls{CCS} \cite{milnerCCS}. Thereby, the components of the system -- for example, clients and servers -- are called processes and are described as algebraic structures. One limitation of the expressiveness of the process algebras mentioned above is the static definition of the processes. A structural modification of the process during a computation is not possible. This is taken into account in the \picalc{} defined by Robin Milner, Joachim Parrow, and David Walker in \cite{milnerParrowWalker}.

The \picalc{} is an extension of \gls{CCS} and adds the facility of expressing changes in the process structure inspired by the work of Uffe Engberg and Mogens Nielsen in \cite{engbertNielsen}. This modification of the process structure is called \emph{mobility}. Thus, with the \picalc{} we can describe and investigate concurrently communicating mobile processes. In the \picalc{}, mobility can be understood as modifications of the \emph{linking structure} of the processes.

A standard example of mobility in the \picalc{} is the connection of mobile phones to base stations. Consider, for example, two persons $P_1$ and $P_2$ carrying mobile phones. At the beginning of the setting both are separately connected to two different static base stations $S_1$ and $S_2$. Thereby, $P_1$ is connected to $S_1$ and $P_2$ is connected to the base station $S_2$. Then, we can think of two possible forms of movement. On the one hand, there is the \emph{physical mobility}. That is, a person has changed its location within the real word, which is not within the focus of the \picalc{}. On the other hand, the communication links between the base stations and the persons has changed, called \emph{virtual mobility}. Consider, for example, $P_1$ is leaving the reception area of $S_1$ while reaching the area of $S_2$. Then a new communication link between $P_1$ and $S_2$ is established and the contact to $S_1$ is disconnected. Thus, $P_1$ and $P_2$ are now both connected to $S_2$ and none of them to $S_1$. Hence, the connections between the persons and the base stations at the beginning of the setting are significantly different to the linking structure at the end. The \picalc{} treats the virtual mobility by passing channels within a communication.

Even though many variations and modifications of Milner's original \picalc{} exist and it has been investigated in many ways, the general focus is mainly on bisimulations between processes to determine their equivalence \cite{pistore}. That is, observing their behaviour and determine whether they behave identically. Thus, proving correctness requires a symmetric equivalence relation with the program and the specification within. Apart from that, it would be interesting to investigate the process' behavior like it is done in \gls{CSP} with its set-theoretical denotation for the investigation whether a process satisfies its specification. That is, we define a condition (the specification) and check if another process (the program) meets this specification by comparing their behavior collected in sets. To reach this, the sequences of events (\emph{traces}) are collected in a set and the correctness is proved by set inclusion. Thereby, there is no need for a symmetry. This approach is often called \emph{refinement}. Even though there is also a relation called simulation which does not need to be symmetrical, ``[$\ldots$] the time has come to unify the two modelling styles [(\gls{CSP} and \gls{CCS})], to enable practicing engineers to exploit a combination of their complementary advantages''\footnote{\cite{hoareWhyCSP}, page 209.} like Tony Hoare wrote in 2006. This statement is also applicable to \gls{CSP} and the \picalc{}.

Hence, the contribution of this thesis is to develop a notion of refinement for \picalc{} processes similar to the approach for \gls{CSP} in \cite{roscoe} to achieve the advantages of this modeling style. Thereby, we primarily concentrate on the \emph{trace semantics} and neglect the \emph{failures} and \emph{failures-divergence semantics}. Furthermore, we investigate the properties of this semantics to show its usability and point out its limitations.

This thesis is divided into five chapters. In \refChap{sec_preliminaries} we briefly introduce sequences and properly investigate the \picalc{} and its operational semantics (the \emph{early transition system} \cite{sangiorgi}). Thereby, a new structural form of processes -- the \emph{extended standard form} -- is introduced and it is shown that every process can be transformed into a structural congruent one in this standard form. We proceed in \refChap{sec_big-step_semantics} by defining a big-step semantics which bases on the early transition system by observing an arbitrary number of external steps within the early semantics. Subsequently, in \refChap{sec_de_sem_trace}, we introduce the trace semantics which is based on the big-step semantics of the previous chapter. Thereby, we investigate its properties and define the refinement based on the trace semantics. Finally, the conclusion in \refChap{sec_conclusion} gives a brief summary of our results and presents ideas for future work.



%oft auch sicht sicht von reaktionen (chemical machine) aber da kann nicht das ganze verhalten usw.
\begin{old}[Milner Zitat]
>>>>>>>>>>>>>>>>>>>>>>>>>>>>>>Zitat Interview Milner: 2003
Even now I am talking with Tony Hoare, who is much more interested in the idea of what it means for a program to meet its specification. We are now trying to reconcile the CCS approach that regards denotations as congruence classes, and the CSP which talks about set-theoretic denotations such as failures and traces and so on, and has very nice ordering relations, so that if the specification is larger than the implementation it means that the implementation is correct. The ordering notion between processes is of course the other main important thing in process calculi. I think we want both; we want the notion of denotations, and perhaps they are congruence classes, and we want the orderings or preorderings representing improvement of, or refinement from, a specification.

> What is the inverse of these refinement relations? I suspect it is related to transactions?

I suppose there is not a symmetry here. If somebody states a specification, then what they are saying is: I want a program to satisfy this and there would be many different programs that satisfy it. To produce any one program that satisfies this is satisfactory. Given that program, you could refine it further and it would of course still satisfy the original specification because that is the refinement ordering. But coarsening it might mean that it then would not satisfy that specification that it originally did, because you would coarsen it in a different direction. So coarsening doesn't have the same pragmatic interest. I think that's a fair way of putting it.


> Yes, but isn't refinement the inverse to that? You have a command that says \"send message from x to y\" and when you refine that, it becomes \"send message from x to a, from a to b, from b to c and from c to y\" ...

I think that's a different kind of refinement. Say you state a specification in logic. You say that all communication must satisfy the following formulae. That allows a whole space of programs to satisfy that. But they will all be talking about communications at a particular level of abstraction. You've got the idea of different levels of abstraction as orthogonal to the question of refinement. Maybe some people will call abstraction refinement, maybe there's confusion ...

>Probably me who confuses things!

But I think everybody is confused about that. Different levels of abstraction is not the same thing as refining from a specification to an implementation. 
\end{old}
%%%%%%%%%%%%%%%%%%%%%%%%%%%%%%%%%%%%%%%%%%%%%% OLD PROPOSAL
\begin{old}[proposal]
In many cases of modern computing it is of interest to describe and model communication, for example, the interaction between components of a computer or -- in particular in our web-based world -- communication across the network. The participating parts often behave in a concurrent manner, hence communication and concurrency are some of the fundamental concepts of computing \cite{milner}.

For the formal specification of such systems process algebras have been widely accepted, especially Robin Milner's \picalc{} described in \cite{milner}. In addition to algebras like the \gls{CCS}, also invented by Robin Milner, and Tony Hoare's \gls{CSP}, the \picalc{} provides the ability to model mobility of processes in an intuitive way, so that the facility of expressing changes in the process structure is given \cite{canal}.

Even though many variations and modifications of Milner's original \picalc{} exist and it is investigated in many ways, the general focus is mainly put on bisimulations between processes to determine their equivalence \cite{pistore}, that is, observing their behaviour and determine whether they behave identically. Apart from that it, would be interesting to distinguish processes by their different behaviour in the degree of nondeterminism and divergence like the refinement notion described by Bill Roscoe in \cite{roscoe} for CSP. 

Thus, the aim of this thesis is to develop a notion of refinement for processes of the \picalc{} and to investigate some of its properties, that is, to construct the possibility to distinguish two \picalc{} processes in the matter of their internal and external behaviour in terms of determinism and divergence. This formerly not reviewed issue is investigated in this theses.

Therefore \refChap{sec_preliminaries} gives an overview of Milner's \picalc{} with its definition and some properties. It also contains some related work \todo[inline]{...}.
\end{old}


	% mainfile: ../Refinement.tex
\chapter{Preliminaries}
\label{sec_preliminaries}
% mainfile: ../../Refinement.tex
At the heart of the refinement of \findex{\picalc{}} processes is the theory of \findex[sequence]{sequences}. Thus, in this chapter, we recall the model of sequences to gain a formal construct to handle ordered elements.% in an intuitive way.

Furthermore, we introduce the \picalc{} and investigate its behavior properly. In particular, we carefully explain the operational semantics of \picalc{} processes, since its peculiarities induce the characteristics of the refinement and its properties. Moreover, we discuss why we choose this particular operational semantics for the following work in this thesis and compare it to other semantics.

The majority of those definitions and notions can, for example, be found in \cite{milner,sangiorgi}.

As mathematical notations, we consider the natural numbers starting with zero ($\N=\set{0,1,2,\ldots}$) and use $\fatsemi$ as the composition of relations. Furthermore, we denote $R^*$ as the reflexive and transitive closure of a relation $R$.


\section{The \texorpdfstring{$\pi$}{pi}-calculus}
\label{sec_pi_calculus}
% mainfile: ../../Refinement.tex
The \findex[\picalc{}|(]{\picalc{}} belongs to the family of \findex[process!algebra]{process algebras} and is used to model communicating and concurrent systems of mobile processes with changing structure. The roots of process algebras can be found in Tony Hoare's \gls{CSP} from the late $70$'s and Milner's \gls{CCS} invented around $1980$ \cite{milnerCCS}. Moreover, the \picalc{} can be seen as a continuation and extension of \gls{CCS} and is inspired by the work of Engberg and Nielsen in \cite{engbertNielsen}, where they extended \gls{CCS} with mobility \cite{milnerParrowWalker}. 

The first published paper treating the \picalc{} is called ``A calculus of mobile processes'' \cite{milnerParrowWalker} written by Robin Milner, Joachim Parrow, and David Walker in the late $80$'s \cite{milner}. While in \gls{CCS} only messages can be sent over a channel during a communication, in this paper they developed the capability to send a channel over a channel. In particular, messages and channels have the same type.%Therein the \picalc{} is defined and it is presented how the calculus can be used to model communicating systems with processes with changing structure; also notions like \findex{strong bisimilarity} and \findex{strong equivalence} are already defined.

Hence, the main advantage of the \picalc{} is its capability to send names over a channel during a communication which can later be used as channels themselves. This provides the ability to change the connection structure of a system.

In this thesis we choose a \picalc{} variant which uses parameterized recursion instead of the replication operator. Thereby, we have no loss in expressiveness, since there is a transformation from recursion to replication and vice versa \cite{milner,sangiorgi}. But since the definition of the examples in the topic of computer science with recursion appeared to be more intuitive to us and recursion is present in \gls{CCS}, we decided to use a variant with recursion rather than replication.%We use this variant at the expense of loosing the possibility of using structural induction \cite{sangiorgi}\todo{instead: fixed point induction?}, because the definition of the examples in the topic of computer science seemed more intuitive to us with recursion.

Furthermore, we define a variant without the match prefix, since ``[m]atch, and especially mismatch, prefixes are seldom useful for describing systems''\footnote{\cite{sangiorgi}, page $13$.}, but inflating the proofs and definitions.

Apart from that, the definition of the syntax as well as the definition of the semantics and the treatment of the \picalc{} refers to \cite{sangiorgi}. Especially the \findex[guarded choice]{guarded choice} -- which means every process in a choice has to perform a prefix action first -- differs from Milner's, Parrow's and Walker's syntax in \cite{milnerParrowWalker}. This restriction is well-accepted in literature \cite{meyer}, since non-guarded choice ``is considered of minor practical importance''\footnote{\cite{meyer}, page $15$.} and ``complicates the theory''\footnote{\cite{sangiorgi}, page $13$.}. However, using guarded choice ``does not delimit the computational expressiveness of the calculus''\footnote{\cite{meyer}, page $15$.}. We slightly differ from the standard definition of the guarded choice. That is, we are not only allowed to have prefix processes in a choice, but we additionally allow that such processes may be wrapped in restrictions. This syntactically extension is needed to transform an arbitrary process into a process in our extended standard form. However, this adaption is of minor influence for the remaining theory.

Finally, we use the \index{\picalc{}!monadic}\findex[monadic \picalc{}]{monadic} \picalc{} -- meaning just a single name can be sent (respectively received) by an output (respectively input) action --, since Milner also gave a transformation from the \index{\picalc{}!polyadic}\findex[polyadic \picalc{}]{polyadic} to the monadic version, which proves their equal expressiveness \cite{milner} and the monadic variant simplifies the handling of processes in this thesis.

\subsection{Syntax}
\label{sec_pi_syntax}
\input{chapters/preliminaries/pi_syntax}

\subsection{Names and substitution}
\label{sec_pi_names_substitution}
\input{chapters/preliminaries/pi_names_substitution}

\subsection{Structural congruence and extended standard form}
\label{sec_pi_struct_cong}
\input{chapters/preliminaries/pi_structural_congruence}

\subsection{Operational semantics}
\label{sec_pi_op_sem}
\input{chapters/preliminaries/pi_operationalSemantics}

\subsection{Simulation and bisimulation}
\label{sec_pi_simulation}
\input{chapters/preliminaries/pi_simulation}


\section{The OZ}
\label{sec_oz}
% mainfile: ../../Refinement.tex
The OZ

\newpage %%%%%%%%%%%%%%% NEWPAGE!





\chapter{Big-step semantics}
\label{sec_big-step_semantics}
% mainfile: ../../Refinement.tex
For retrieving the whole behavior of a process, we investigate the transition system produced by the rules of \refFig{fig_ts_early}. In particular, we follow the way of the transitions and collect their labels. Thus, some kind of transition-chains -- so-called \findex[trace]{traces} -- are produced. In this chapter we define a big-step semantics by stringing together the steps of the early semantics as well as the formal meaning of traces. The ideas of the definition of the big-steps and traces base upon the trace model for \gls{CSP} introduced by Hoare in \cite{hoareTraceOriginal} and \cite{hoare}.%\todo{or \cite{roscoe}. justify}

Additionally, in this chapter some of the definitions and properties explained in \refSec{sec_pi_calculus} are lifted to traces, since the early transition system is at the heart of the big-step semantics.

\section{Definition}
\label{sec_big-step_def}
% mainfile: ../../Refinement.tex
Since we would like to follow a chain of transitions and its labels through the transition system, we define a formal construct which collects actions in an ordered way. Moreover, we abstract from the internal behavior. Thus, we are only collecting visible actions.

\begin{definition}[Traces]
\label{def_traces}\index{trace}
	We define a set containing all sequences of external actions by
		\[\tr \define \seqset{\actions\setminus\set{\tau}}\]
	and call each of its elements a \findex{trace}. In particular, we call $\eseq\in\tr$ the \findex{empty trace}.
\end{definition}

Thus, with a trace, it is possible to describe the progression of external actions. In the big-step semantics this is combined with the operational semantics to easily notate by which external actions a process can evolve to another.

%%%%%%%%%%%%%%%%%%%%%%%%%%%%%%%%%%%%%%%%%%%%%%%%%%%%%%%%%%% START OLD DEFINITION %%%%%%%%%%%%%%%%%%%%%%%%%%%%%%%%%%%%%%%%%%%%%%%%%%%%%%%%%%%%%%%%
\begin{old}{old definition}
\begin{definition}[Big-step semantics]
\label{def_bigstep_semantics}
	The \findex{bigstep semantics} $\set[s\in\tr]{\bigstep{s}}\subseteq\procs\times\procs$ is inductively defined. For processes $P,Q\in\procs$:
		\[P\bigstep{\eseq}Q, \falls \ec{P}(\tautrans)^*\ec{Q}\]
	where $(\tautrans)^*$ is the reflexive and transitive closure of $\tautrans$.
		\begin{align*}
			P\bigstep{\seq{a}}Q,\falls& P',Q'\in\procs\text{ exists such that} \\
			&P\bigstep{\eseq}P'\;and\;\ec{P'}\transs{a}\ec{Q'}\;and\;Q'\bigstep{\eseq}Q
		\end{align*}
	with $a\in\actions\setminus\set{\tau}$.
		\[P\bigstep{\seqconc{s}{t}}Q,\falls P\bigstep{s}\fatsemi\bigstep{t}Q\]
	with $s,t\in\tr$ and $\fatsemi$ the composition of relations.
\end{definition}
\end{old}
%%%%%%%%%%%%%%%%%%%%%%%%%%%%%%%%%%%%%%%%%%%%%%%%%%%%%%%%%%% END OLD DEFINITION %%%%%%%%%%%%%%%%%%%%%%%%%%%%%%%%%%%%%%%%%%%%%%%%%%%%%%%%%%%%%%%%

\begin{definition}[Big-step semantics]
\label{def_bigstep_semantics}
	We inductively define the \findex{big-step semantics} $\set[s\in\tr]{\bigstep{s}}\subseteq\procs_\alpha\times\procs_\alpha$ for processes $P,Q\in\procs$ by
		\[\ec{P}\bigstep{\eseq}\ec{Q},\;\gdw\;\ec{P}(\tautrans)^*\ec{Q}.\]
	As a shorter notion it is also possible to write $\ec{P}\bigstep{}\ec{Q}$.
		\begin{align*}
			\ec{P}\bigstep{\seq{\alpha}}\ec{Q},\;\gdw\;\ec{P}\bigstep{\eseq}\fatsemi\transs{\alpha}\fatsemi\bigstep{\eseq}\ec{Q}
		\end{align*}
	with $\alpha\in\actions\setminus\set{\tau}$.
		\[\ec{P}\bigstep{\seqconc{s}{t}}\ec{Q},\;\gdw\;\ec{P}\bigstep{s}\fatsemi\bigstep{t}\ec{Q}\]
	with $s,t\in\tr$.

	Furthermore, we call $\bigstep{\eseq{}}$ an \findex[invisible big-step]{invisible big-step} and $\bigstep{t}$ for a trace $t\in\tr$ a \findex{visible big-step}.
\end{definition}

Thus, we now have the possibility to write $\ec{\procpar{\inp{a}{x}.\out{x}{z}}{\out{a}{y}}}\bigstep{\seq{\inpa{a}{b},\out{a}{y},\out{b}{z}}}\ec{\procpar{\proczero}{\proczero}}$, since $\ec{\procpar{\inp{a}{x}.\out{x}{z}}{\out{a}{y}}}\intrans{a}{b}\ec{\procpar{\out{b}{z}}{\out{a}{y}}}\outtrans{a}{y}\ec{\procpar{\out{b}{z}}{\proczero}}\outtrans{b}{z}\ec{\procpar{\proczero}{\proczero}}$. Furthermore, we abbreviate $\ec{\procpar{\inp{a}{x}.\out{x}{z}}{\out{a}{y}}}\bigstep{\seq{\inpa{a}{b},\out{a}{y}}}\ec{\procpar{\out{b}{z}}{\proczero}}$ and $\ec{\procpar{\out{b}{z}}{\proczero}}\bigstep{\seq{\out{b}{z}}}\ec{\procpar{\proczero}{\proczero}}$ by $\ec{\procpar{\inp{a}{x}.\out{x}{z}}{\out{a}{y}}}\bigstep{\seq{\inpa{a}{b},\out{a}{y}}}\ec{\procpar{\out{b}{z}}{\proczero}}\bigstep{\seq{\out{b}{z}}}\ec{\procpar{\proczero}{\proczero}}$.

Note that $\ec{P}\bigstep{t}\ec{P_1}\transs{\alpha}\ec{Q}$ for some $\alpha\in\actions$ and $\ec{P}\bigstep{t}\ec{P_1'}$ in general does not imply that $\ec{P_1'}$ must have a possibility for a transition labeled with $\alpha$. Consider, for example, $P\define\procchoice{\out{a}{x}}{\out{a}{x}.\out{b}{y}}$. Then $\ec{P}\bigstep{\seq{\outa{a}{x}}}\ec{\out{b}{y}}\outtrans{b}{y}\ec{\proczero}$ and $\ec{P}\bigstep{\seq{\outa{a}{x}}}\ec{\proczero}$.


\section{Names and binding}
\label{sec_big-step_name_binding}
% mainfile: ../../Refinement.tex
Since we already talked about objects and subjects of actions, it is useful to define functions which calculate them from the actions and especially from traces.

\begin{definition}[Subjects and objects]
\label{def_sub_obj}
	We define $\subF: \actions\setminus\set{\tau} \rightarrow \names$ as the function which collects the \findex[subject]{subjects} and $\objF: \actions\setminus\set{\tau} \rightarrow \names$ for the \findex[object]{objects} of an action such that
	\begin{equation*}
		\begin{aligned}
			\begin{array}{lcl}
				\sub{\out{a}{x}} & \define & a \\
				\sub{\bout{a}{x}} & \define & a \\
				\sub{\inpa{a}{x}} & \define & a 
			\end{array}
		\end{aligned}\quad\quad\quad\quad\quad\quad
		\begin{aligned}
		\begin{array}{lcl}
			\obj{\out{a}{x}} & \define & x \\
			\obj{\bout{a}{x}} & \define & x \\
			\obj{\inpa{a}{x}} & \define & x 
		\end{array}
	\end{aligned}
\end{equation*}
holds. With those functions, we define $\subF: \tr \rightarrow \pom{\names}$ with
	\[\sub{s} \define \set[\exists\alpha\in{}s: x=\sub{\alpha}]{x\in\names},\]
which collects the subjects of a trace and $\objF: \tr \rightarrow \pom{\names}$ with
\[\obj{s} \define \set[\exists\alpha\in{}s: x=\obj{\alpha}]{x\in\names},\]
which collects the objects.
\end{definition}

Furthermore, we %already know what the meaning of free and bound names of actions as well as of processes is, it is consistent to
define the \findex[free name!trace]{free} and \findex[bound name!trace]{bound names} of a trace. This definition intuitively fits to the definition of free and bound names of processes in \refSec{sec_pi_names_substitution}.

\begin{definition}[Bound and free names of a trace]
\label{def_bound_free_names_traces}
The \findex[bound name!trace]{bound names} of a trace $t\in\tr$ are defined by
\[\bn{t}\define\bigcup_{\alpha\in{}t}\bn{\alpha}.\]
Thus, all names are collected, which occur as an object of a bound output action within the trace. Furthermore, with
\[\fn{t}\define{}\set[n\nin\bn{t}]{n\in\n{t}}\]
we define the \findex[free name!trace]{free names} of a trace $t\in\tr$ by the set of all the other names, which do not occur bound in $t$. Once more, the \findex[name!trace]{names} of a transition $t\in\tr$ are denoted by $
\n{t}\define{}\fn{t}\cup\bn{t}$.
%\end{align*}
%	\begin{equation*}	
%		\begin{aligned}			
%			\begin{array}{lcl}
%				\fn{\eseq} & \define & \emptyset \\
%				\fn{\seqconc{\seq{\out{a}{x}}}{t}} & \define & \fn{t}\cup\set{a,x} \\
%				\fn{\seqconc{\seq{\inp{a}{x}}}{t}} & \define & \left(\fn{t}\cup\set{a}\right)\setminus\set{x} \\
%				\fn{\seqconc{\seq{\bout{a}{x}}}{t}} & \define & \left(\fn{t}\cup\set{a}\right)\setminus\set{x} 
%			\end{array}
%		\end{aligned}\quad
%		\begin{aligned}
%			\begin{array}{lcl}
%				\bn{\eseq} & \define & \emptyset \\
%				\bn{\seqconc{\seq{\out{a}{x}}}{t}} & \define & \bn{t} \\
%				\bn{\seqconc{\seq{\inp{a}{x}}}{t}} & \define & \bn{t}\cup\set{x} \\
%				\bn{\seqconc{\seq{\bout{a}{x}}}{t}} & \define &\bn{t}\cup\set{x} 
%			\end{array}
%		\end{aligned}
%	\end{equation*}
\end{definition}

Like the connection between the operational semantics and the free names of its processes, there is also a connection between traces and the free names of its components.

\begin{lemma}[Big-step semantics and free names]
\label{lem_bigstep_fn}
Given processes $P,Q\in\procs$ and a trace $t\in\tr$, then
\[\ec{P}\bigstep{t}\ec{Q} \text{ implies } \fn{Q}\subseteq\fn{P}\cup\n{t}\]
holds.
\end{lemma}
\begin{prf}
Let $P_1,P_n\in\procs$ and $t\in\tr$ with $\ec{P_1}\bigstep{t}\ec{P_n}$ and $\len{t}=n-1$ for some $n\in\N$ . Hence, there are processes $P_2,\ldots,P_{n-1}\in\procs$ such that $\ec{P_1}\bigstep{\seq{t_1}}\ec{P_2}\bigstep{\seq{t_2}}\cdots\bigstep{\seq{t_{n-2}}}\ec{P_{n-1}}\bigstep{\seq{t_{n-1}}}\ec{P_n}$. For every $\tau$ step within the big-steps we know from \refLem{lem_trans_fn} Equation~(\ref{eq_trans_fn_tau}) that the free names of the resulting process are a subset of the free names of the starting one. The same lemma also yields for every $i\in\set{1,\ldots,n-1}$ that $\fn{P_{i+1}}\subseteq{}\fn{P_i}\cup\n{t_{i}}$. Thus, $\fn{P_n}\subseteq{}\fn{P_{n-1}}\cup\n{t_{n-1}}\subseteq\left(\fn{P_{n-2}}\cup\n{t_{n-2}}\right)\cup\n{t_{n-1}} \subseteq \cdots \subseteq\fn{P_1}\cup\n{t}$, since $\n{t}=\bigcup_{i\in{}\set{1,\ldots,n-1}}\n{t_i}$ holds.
\end{prf}

Note that in particular the inclusion is an over-approximation. Consider, for instance, $P\define\procpar{\inp{a}{b}.\out{c}{d}}{\out{x}{y}.\inp{z}{w}}$. Then, $\ec{P}\bigstep{\seq{\inpa{a}{b},\outa{x}{y}}}\ec{Q}$ with $Q\define{}\procpar{\out{c}{d}}{\inp{z}{w}}$ and $\fn{Q}=\set{c,d,z}\subseteq\set{a,c,d,x,y,z}\cup\set{a, b, x, y}=\fn{P}\cup\n{t}$. Moreover, $\ec{P}\bigstep{\seq{\inpa{a}{r}, \outa{c}{d}, \outa{x}{y}, \inpa{z}{s}}}\ec{\procpar{\proczero}{\proczero}}$ and so $\emptyset\subseteq\set{a,c,d,x,y,z}\cup\set{a, r, c,d,x,y,z,s}$.

Furthermore, it is helpful -- especially for \refLem{lem_from_resset_to_res} -- to define a function which replaces every occurrence of an output action in a trace by a bound output action for a given name.

\begin{definition}[Binding]
\label{def_binding}
	We define the \findex{binding function} $\bindF: \names\times\actions \rightarrow \actions$, with
	\begin{align*}
		\begin{array}{lcl}
			\bind{a}{\tau} & \define & \tau \\
			\bind{a}{\bout{b}{x}} & \define & \bout{b}{x} \\
			\bind{a}{\inpa{b}{x}} & \define & \inpa{b}{x} \\
			\bind{a}{\out{b}{x}} & \define & \left\{\begin{array}{ll}
									\bout{b}{x} & \falls\; a=x \\
									\out{b}{x} & \text{else}
								\end{array} \right.
		\end{array}
	\end{align*}
	such that it only replaces an output action with its corresponding bound output action, if the given name is the object of the action.

Furthermore, for traces $\bindF: \names\times\tr \rightarrow \tr$ is defined as
\begin{align*}
\bind{a}{s} \define \seq{s'_1,s'_2,\dots} \text{ with }  s'_i =\left\{
								\begin{array}{cll}
									 \bind{a}{s_i} & \text{if } \nexists{}j\in\N\colon{}&j<i \wedge a\in\obj{s_j} \\
													& & \wedge{}s_j\in\outA\cup\boutA \\
									s_i & \text{else}&
								\end{array}\right.
\end{align*}
for all $i\in\set{1,2,\dots}$. Moreover, we would like to bind a name in a whole set of traces such that $\bindF: \names\times\pom{\tr} \rightarrow \pom{\tr}$ is defined with
\[\bind{a}{X} \define \set[\exists s'\in{}X: s=\bind{a}{s'}]{s\in\tr}\]
as a function which binds a name to a whole trace set by applying the binding function to every trace within the set.
\end{definition}

Note that just the first occurrence of the name as an object of an output action is replaced by the binding function. This is because by \refConv{conv_uni_bn_traces} we will stipulate that every bound name within a trace is unique.

%%%%%%%%%%%%%%%%%%%%%%%%%%%%%%%%%%%%%%%%% OLD: wrong explanation %%%%%%%%%%%%%%%%%%%%%%%%%%%%%%%%%%%%%%%%%%%%%%%%%%%%%
\begin{old}{incorrect explanation why there can be at least one occurrence of a bound output}
there can be at most one occurrence of an output action in a trace which have to be replaced. This is because there is just one rule in the definition of the operational semantics (\refDef{def_early_trans_system}) which produces a bound output action and thereby the restriction operator is afterwards omitted and no bound output action, without a restriction. With \refConv{conv_uni_bn}, we know that the given name cannot occur twice bounded in a process, so another bound output action for the same name cannot exist.
\end{old}
%%%%%%%%%%%%%%%%%%%%%%%%%%%%%%%%%%%%%%%%% OLD: end wrong explanation %%%%%%%%%%%%%%%%%%%%%%%%%%%%%%%%%%%%%%%%%%%%%%%%%%%%%


\section{Substitution}
\label{sec_big-step_substitution}
% mainfile: ../../Refinement.tex
%%%%%%%%%%%%%%%%%%%%%%%%%%%%%%%%%%%%%%%%%%%%%%%%%% START NOT NEEDED INDUCTIVE DEFINITION %%%%%%%%%%%%%%%%%%%%%%%%%%%%%%%%%%%%
\begin{old}{not needed inductive definition}
\begin{definition}[Substitution on traces]
\label{def_substitution_traces}
The \index{substitution!trace}\findex[application substitution!trace]{application} of a substitution $\sigma$ to a trace $t\in\tr$ is written $t\sigma\in\tr$ and is inductively defined
	\begin{align*}
		\begin{array}{lcl}
			\eseq{}\sigma &\define& \eseq \\
			\left(\seqconc{\seq{\out{a}{x}}}{t}\right)\sigma &\define& \seqconc{\seq{\out{\sigma(a)}{\sigma(x)}}}{(t\sigma)} \\
			\left(\seqconc{\seq{\inp{a}{x}}}{t}\right)\sigma &\define& \seqconc{\seq{\inp{\sigma(a)}{p}}}{((t\subs{p}{x})\sigma)} \; for \; p\nin\fn{t}\cup\set{a}\cup\n{\sigma}\\
			\left(\seqconc{\seq{\bout{a}{x}}}{t}\right)\sigma &\define& \seqconc{\seq{\bout{\sigma(a)}{p}}}{((t\subs{p}{x})\sigma)}\; for \; p\nin\fn{t}\cup\set{a}\cup\n{\sigma}		\end{array}
	\end{align*}
with $a,x,p\in\names$ and $t\in\tr$. For a set of traces $X\subseteq\tr$ we define the application of the substitution $X\sigma\define\set[t\in{}X]{t\sigma}$ as the application of the substitution to all elements of the set.
\end{definition}
\end{old}
%%%%%%%%%%%%%%%%%%%%%%%%%%%%%%%%%%%%%%%%%%%%%%%%%% END NOT NEEDED INDUCTIVE DEFINITION %%%%%%%%%%%%%%%%%%%%%%%%%%%%%%%%%%%%

Since we can replace names within processes and actions, we consequently define the application of a substitution on traces.

\begin{definition}[Substitution on traces]
\label{def_substitution_traces}
The \index{substitution!trace}\findex[application substitution!trace]{application} of a substitution $\sigma$ to a trace $t\in\tr$ with $\n{\sigma}\cap\bn{t}=\emptyset$ is written $t\sigma\in\tr$ and is defined as the application of the substitution to every action: $t\sigma\define\set[(i,\alpha)\in{}t]{(i,\substF(\alpha))}$. For a set of traces $X\subseteq\tr$ we define the application of the substitution $X\sigma\define\set[t\in{}X]{t\sigma}$ as the application of the substitution to all elements of the set.
\end{definition}

%%%%%%%%%%%%%%%%%%%%%%%%%%%%%%%%%%%%%%%%%%%%%%%%%% START NOT NEEDED DEFINITION %%%%%%%%%%%%%%%%%%%%%%%%%%%%%%%%%%%%
\begin{old}{not needed definition of a bound substitution}
\todo[inline]{In some cases, we also need to replace the bound names of an action respectively trace. Therefore, we define a \findex{bound substitution} function.}
\begin{definition}[Bound substitution]
\label{def_bound_substitution}
As $\bnsubstF{}:\pom{\names^2}\times\tr\rightarrow\tr$ with
\begin{align*}
\left(\subs{x}{y},t\right)\mapsto\left\{ \begin{array}{ll}
					t\subs{x}{y} & \text{if } x\nin\left(\fn{t}\cup\bn{t}\right)\setminus\set{y} \\
					t 		& \text{else}
				\end{array}\right.
\end{align*}
we define the \findex{bound substitution} function and lift it with $\bnsubstF{}:\pom{\names^2}\times\pom{\tr}\rightarrow\pom{\tr}$ and
\[ \left(\subs{x}{y}, T\right)\mapsto \set[\exists{}t\in{}T: t'=\bnsubst{x}{y}{t}]{t'\in\tr}\]
on traces.
\end{definition}
\todo[inline]{The bounding substitution replaces every occurrence of a given name, apart from the name which should be replaced.}
\end{old}
%%%%%%%%%%%%%%%%%%%%%%%%%%%%%%%%%%%%%%%%%%%%%%%%%% END NOT NEEDED DEFINITION %%%%%%%%%%%%%%%%%%%%%%%%%%%%%%%%%%%%
%%%%%%%%%%%%%%%%%%%%%%%%%%%%%%%%%%%%%%%%%%%%%%%%%% START NOT NEEDED DEFINITION %%%%%%%%%%%%%%%%%%%%%%%%%%%%%%%%%%%%
\begin{old}{not needed definition substitution on transitions}
\begin{definition}[Substitution on transitions]
\label{def_substitution_transitions}
The \index{substitution!transition}\findex[application substitution!transition]{application} of a substitution $\sigma$ to a transition $\ec{P}\transs{\alpha}\ec{Q}$ for $P,Q\in\procs$ and $\alpha\in\actions$ is written $\left(\ec{P}\transs{\alpha}\ec{Q}\right)\sigma$ and is defined
	\begin{align*}
		\left(\ec{P}\transs{\alpha}\ec{Q}\right)\sigma\define
					\left\{\begin{array}{ll}
							\ec{P}\tautrans\ec{Q} & \falls\;\alpha=\tau\\
							\ec{P}\outtrans{\sigma(a)}{\sigma(x)}\ec{Q} & \falls\;\alpha=\out{a}{x}\\
							\ec{P}\intrans{\sigma(a)}{p}\ec{Q\subs{p}{x}}& \falls\;\alpha=\inp{a}{x}\\
							\ec{P}\bouttrans{\sigma(a)}{p}\ec{Q\subs{p}{x}}& \falls\;\alpha=\bout{a}{x}
						\end{array}
		\right.
	\end{align*}
with $p\nin\left(\fn{P}\cup\n{\sigma}\right)$.
\end{definition}
\end{old}
%%%%%%%%%%%%%%%%%%%%%%%%%%%%%%%%%%%%%%%%%%%%%%%%%% END NOT NEEDED DEFINITION %%%%%%%%%%%%%%%%%%%%%%%%%%%%%%%%%%%

With \refLem{lem_subst_trans_partI} we showed a relation between $\tau$ transitions and substitutions. In the following, we lift this lemma to the big-step semantics. Thus, we know if there is an invisible big-step between two processes, there is also an invisible big-step between the processes where a substitution has been applied to.

\begin{lemma}[Substitution on big-step semantics (Part I)]
\label{lem_subst_bigstep_partI}
For all processes $P,Q\in\procs$ and substitutions $\sigma$,
\[\text{if } \ec{P}\bigstep{\eseq}\ec{Q} \text{ then } \ec{P\sigma}\bigstep{\eseq}\ec{Q\sigma}\]
holds.
\end{lemma}
\begin{prf}
Let $P,Q\in\procs$ and $\sigma$ a substitution, then
\begin{align*}
	\begin{array}{lcl}
		\ec{P}\bigstep{\eseq}\ec{Q} & \stackrel{\{\text{\refDef{def_bigstep_semantics}\}}}{\text{impl.}} & \ec{P}\left(\tautrans\right)^*\ec{Q} \\
			& \text{impl.} & \exists{}Q_1,\ldots,Q_n\in\procs:\ec{P}\tautrans\ec{Q_1}\wedge\ec{Q_1}\tautrans\ec{Q_2} \\
					&&\quad\quad\quad\quad\quad\quad\quad\quad\quad \wedge \cdots\wedge\ec{Q_n}\tautrans\ec{Q} \\		
			& \stackrel{\{\text{Lem.~\ref{lem_subst_trans_partI}\}}}{\text{impl.}} & \exists{}Q_1,\ldots,Q_n\in\procs:\ec{P\sigma}\transs{\sigma(\tau)}\ec{Q_1\sigma}\wedge\ec{Q_1\sigma}\transs{\sigma(\tau)}\ec{Q_2\sigma} \\
					&&\quad\quad\quad\quad\quad\quad\quad\quad\quad \wedge \cdots\wedge\ec{Q_n\sigma}\transs{\sigma(\tau)}\ec{Q\sigma}\\
			& \text{impl.} & \ec{P\sigma}\left(\tautrans\right)^*\ec{Q\sigma} \\
			& \stackrel{\{\text{Def.~\ref{def_bigstep_semantics}\}}}{\text{impl.}} & \ec{P\sigma}\bigstep{\eseq}\ec{Q\sigma}
	\end{array}
\end{align*}
holds.
\end{prf}

Since \refLem{lem_subst_trans_partI} is at the heart of this proof, we see that the converse of this lemma does not hold. Consider, for example, $P\define\procpar{\inp{a}{x}}{\out{b}{y}}$ and $\sigma\define\subs{a}{b}$. Then $\ec{P\sigma}\bigstep{}\ec{Q\sigma}$, with $Q\define\procpar{\proczero}{\proczero}$ holds, but in $\ec{P}$ no $\tau$ transition is possible. However, there is also a restricted version of the converse of this lemma presented in \refLem{lem_subst_bigstep_partIII}.

%We lift the idea of \refLem{lem_subst_bigstep_partI} to traces with visible actions, we need a restriction to the trace for which the big-step between two processes leads to a big-step with an application of a substitution to all of the components. Therefore, we define traces where the bound names within are unique.

We lift the idea of \refLem{lem_subst_bigstep_partI} to traces with visible actions such that the big-step between two processes leads to a big-step with an application of a substitution to all of its components. Thereby, the lifting is limited to substitutions different to the bound names of a trace. Later, we will see with \refConv{conv_uni_bn_traces} that this restriction is no longer harmful.

\begin{lemma}[Substitution on big-step semantics (Part II)]
\label{lem_subst_bigstep_partII}
Given process"-es $P,Q\in\procs$, a trace $s\in\tr$ and a substitution $\sigma$ with $\n{\sigma}\cap\bn{s}=\emptyset$,
\[\text{if } \ec{P}\bigstep{s}\ec{Q} \text{ then } \ec{P\sigma}\bigstep{s\sigma}\ec{Q\sigma}\]
holds.
\end{lemma}
\begin{prf}
Let $P,Q\in\procs$, $s\in\tr$ and $\sigma$ a substitution with $\n{\sigma}\cap\bn{s}=\emptyset$. We proceed by induction over the length of $s$.
\begin{description}
\item[Base case $s=\eseq$:] \refLem{lem_subst_bigstep_partI} directly yields the proof of this case.

\item[Base case $s=\seq{\alpha}$:] Let $\alpha\in\actions\setminus\set{\tau}$ and $\ec{P}\bigstep{s}\ec{Q}$ hold. So we know from the definition of the big-step semantics (\refDef{def_bigstep_semantics}) that there have to be processes $Q_i$ and $Q_j$ such that $\ec{P}\bigstep{}\ec{Q_i}\transs{\alpha}\ec{Q_j}\bigstep{}\ec{Q}$. With \refLem{lem_subst_bigstep_partI} we know that $\ec{P\sigma}\bigstep{}\ec{Q_i\sigma}$ and $\ec{Q_j\sigma}\bigstep{}\ec{Q\sigma}$ hold. Furthermore, \refLem{lem_subst_trans_partI} yields that with $\ec{Q_i}\transs{\alpha}\ec{Q_j}$, also $\ec{Q_i\sigma}\transs{\sigma(\alpha)}\ec{Q_j\sigma}$ holds since $\n{\sigma}\cap\bn{s}=\emptyset$. So we know $\ec{P\sigma}\bigstep{\seq{\sigma(\alpha)}}\ec{Q\sigma}$, hence $\ec{P\sigma}\bigstep{s\sigma}\ec{Q\sigma}$ holds.

\item[Induction hypothesis:] For a number $n\in\N$ and all traces $s\in\tr$, with $n\leq\#(s)$ holds: For all $P,Q\in\procs$ and $\sigma$ substitution with $\n{\sigma}\cap\bn{s}=\emptyset$ if $\ec{P}\bigstep{s}\ec{Q}$ then $\ec{P\sigma}\bigstep{s\sigma}\ec{Q\sigma}$. 

\item[Inductive step $n\mapsto n+1$:] Let $s'\in\tr$ with $\#(s')=n+1$ and $\ec{P}\bigstep{s'}\ec{Q}$. So we know that a trace $s\in\tr$ and an action $\alpha\in\actions\setminus\set{\tau}$ exist such that $s'=\seqconc{s}{\seq{\alpha}}$. Then
\begin{align*}
	\begin{array}{lcl}
		\ec{P}\bigstep{\seqconc{s}{\seq{\alpha}}}\ec{Q} & \stackrel{\{\text{\refDef{def_bigstep_semantics}\}}}{\text{impl.}} & \ec{P}\bigstep{s}\fatsemi\bigstep{\seq{\alpha}}\ec{Q} \\
			& \text{impl.} & \exists{}Q'\in\procs:\ec{P}\bigstep{s}\ec{Q'}\wedge\ec{Q'}\bigstep{\seq{\alpha}}\ec{Q}\\
			& \stackrel{\{\text{IH\}}}{\text{impl.}} & \exists{}Q'\in\procs:\ec{P\sigma}\bigstep{s\sigma}\ec{Q'\sigma}\wedge\ec{Q'\sigma}\bigstep{\seq{\alpha}\sigma}\ec{Q\sigma}\\
			& \text{impl.} & \ec{P\sigma}\bigstep{s\sigma}\fatsemi\bigstep{\seq{\alpha}\sigma}\ec{Q\sigma}\\
			& \text{impl.} & \ec{P\sigma}\bigstep{\seqconc{s\sigma}{\seq{\alpha}\sigma}}\ec{Q\sigma}\\
			& \text{impl.} & \ec{P\sigma}\bigstep{\left(\seqconc{s}{\seq{\alpha}}\right)\sigma}\ec{Q\sigma}\\
			& \text{impl.} & \ec{P\sigma}\bigstep{s'\sigma}\ec{Q\sigma}
	\end{array}
\end{align*}
holds.
\end{description}
Thus, \refLem{lem_subst_bigstep_partII} holds.
\end{prf}

With another small example, we can see that the converse of this lemma does not hold either. Consider, for example, $P\define{}\procpar{\inp{a}{x}.\out{x}{c}}{\out{b}{y}.\out{c}{d}}$ and $\sigma\define{}\subs{a}{b}$. Then $\ec{P\sigma}=\ec{\procpar{\inp{a}{x}.\out{x}{c}}{\out{a}{y}.\out{c}{d}}}\tautrans\ec{\procpar{\out{y}{c}}{\out{c}{d}}}\outtrans{y}{c}\ec{\procpar{\proczero}{\out{c}{d}}}$, so $\ec{P\sigma}\bigstep{\seq{\out{y}{c}}}\ec{\procpar{\proczero}{\out{c}{d}}}$ holds. But there is no trace $s\in\tr$ with $s\sigma=\seq{\out{y}{c}}$ and no process $Q\in\procs$ with $\ec{Q\sigma}=\ec{\procpar{\proczero}{\out{c}{d}}}$ such that $\ec{P}\bigstep{s}\ec{Q}$, because every trace starting in $\ec{P}$ has right at the front an input action or $\out{b}{y}$. And since $y\nin\cosupp{\sigma}$, we know that for no trace $s$ starting in $\ec{P}$ can $s\sigma=\seq{\out{y}{c}}$ hold.

Furthermore, consider, for example, $P\define\inp{y}{a}$ and $\sigma\define\subs{b}{x}$, then $\ec{P\sigma}=\ec{P}\bigstep{\seq{\inpa{y}{x}}}\ec{\proczero}$. Thus, the converse would yield that there must be a trace $s\in\tr$ and a process $Q$ with $\ec{P}\bigstep{s}\ec{Q}$ and $s\subs{b}{x}=\seq{\inpa{y}{x}}$. But we cannot find a trace $s$ such that we can apply a substitution to $s$ which replaces the name $x$ and results in a trace where the name $x$ occurs free.

But as in the operational semantics case there are restricted converses such that we gain a connection between big-steps with and without the application of a substitution.

\begin{lemma}[Substitution on big-step semantics (Part III)]
\label{lem_subst_bigstep_partIII}
Given $a,b\in\names$, $P\in\procs$ and 
\begin{itemize}
\item[(I)] $\sigma=\subs{a}{b}$ a substitution with $a\nin\fn{P}$ or
\item[(II)] $\sigma=\transp{a}{b}$ a transposition,
\end{itemize} then
\[\ec{P\sigma}\bigstep{}\ec{Q'} \text{ implies } \exists Q\in\procs:\ec{P}\bigstep{}\ec{Q}\]
with $\ec{Q\transp{a}{b}}=\ec{Q'}$ holds.
\end{lemma}
\begin{prf}
Let $a,b\in\names$, $P\in\procs$ and $\sigma=\subs{a}{b}$ substitution with $a\nin\fn{P}$ or $\sigma=\transp{a}{b}$ transposition. Furthermore, there is a $Q'\in\procs$ such that $\ec{P\sigma}\bigstep{}\ec{Q'}$. We proceed by induction over the number $n\in\N$ of $\tau$ steps within $\bigstep{}$.
\begin{description}
\item[Base case $n=0$:] Thus, $\ec{P\sigma}=\ec{Q'}$. Chose $Q\define{}P$, so we know $\ec{P}\bigstep{}\ec{Q}$ and $\ec{Q\transp{a}{b}}=\ec{P\transp{a}{b}}=\ec{P\sigma}=\ec{Q'}$ holds.

\item[Base case $n=1$:] Thus, $\ec{P\sigma}\tautrans{}\ec{Q'}$. \refCor{cor_subst_trans_partIII} respectively \refCor{cor_subst_trans_partIV} yields that there is an action $\alpha\in\actions$ and a process $Q\in\procs$ such that $\ec{P}\transs{\alpha}\ec{Q}$ and $\transp{a}{b}(\alpha)=\tau$ and $\ec{Q\transp{a}{b}}=\ec{Q'}$. Hence, $\alpha=\tau$ and so $\ec{P}\bigstep{}\ec{Q}$.

\item[Induction hypothesis:] For all $a,b\in\names$, $P\in\procs$ and $\sigma=\subs{a}{b}$ a substitution with $a\nin\fn{P}$ or $\sigma=\transp{a}{b}$ transposition and $\ec{P\sigma}\bigstep{}\ec{Q'}$ with some number of $\tau$ transitions $n\in\N$, exists a process $Q\in\procs$ with $\ec{P}\bigstep{}\ec{Q}$ and $\ec{Q\transp{a}{b}}=\ec{Q'}$.

\item[Induction step:] Let $\ec{P\sigma}\bigstep{}\ec{Q'}$ with a number of $\tau$ steps of $n+1$. Thus, we know there is a process $Q_1'\in\procs$ such that $\ec{P\sigma}\bigstep{}\ec{Q_1'}$ in a number of $\tau$ steps of $n$ and $\ec{Q_1'}\tautrans{}\ec{Q'}$. From the induction hypothesis we know that there is a process $Q_1\in\procs$ such that $\ec{P}\bigstep{}\ec{Q_1}$ with $\ec{Q_1\transp{a}{b}}=\ec{Q_1'}$. So we know $\ec{Q_1\transp{a}{b}}\tautrans{}\ec{Q'}$ and with \refCor{cor_subst_trans_partIII} respectively \refCor{cor_subst_trans_partIV} we know there is an action $\alpha\in\actions$ and a process $Q\in\procs$ such that $\ec{Q_1}\transs{\alpha}\ec{Q}$ and $\transp{a}{b}(\alpha)=\tau$ and $\ec{Q\transp{a}{b}}=\ec{Q'}$. Thus, $\alpha=\tau$ and due to that $\ec{P}\bigstep{}\ec{Q}$.
\end{description}
Thus, \refLem{lem_subst_bigstep_partIII} holds.
\end{prf}

With the help of \refCor{cor_subst_trans_partIII} we extend the previous lemma for big-steps with visible actions.

\begin{lemma}[Substitution on big-step semantics (Part IV)]
\label{lem_subst_bigstep_partIV}
Given names $a,b\in\names$, a process $P\in\procs$, a trace $s'\in\tr$, and 
\begin{itemize}
\item[(I)] $\sigma=\subs{a}{b}$ a substitution with $a\nin\fn{P}$ or
\item[(II)] $\sigma=\transp{a}{b}$ a transposition
\end{itemize}
with $\n{\sigma}\cap\bn{s'}=\emptyset$, then 
\[\ec{P\sigma}\bigstep{s'}\ec{Q'} \text{ implies } \exists Q\in\procs,s\in\tr:\ec{P}\bigstep{s}\ec{Q}\]
with $s\transp{a}{b}=s'$ and $\ec{Q\transp{a}{b}}=\ec{Q'}$.
\end{lemma}
\begin{prf}
Let $a,b\in\names$, $P\in\procs$, $s'\in\tr$, and $\sigma=\subs{a}{b}$ a substitution with $a\nin\fn{P}$ or $\sigma=\transp{a}{b}$ transposition with $\n{\sigma}\cap\bn{s'}=\emptyset$ and $\ec{P\sigma}\bigstep{s'}\ec{Q'}$. We proceed by induction over the length of $s'$.
\begin{description}
\item[Base case $\len{s'}=0$:] Thus, $s'=\eseq{}$ holds. Hence, in this case \refLem{lem_subst_bigstep_partIV} directly follows from \refLem{lem_subst_bigstep_partIII}.

\item[Base case $\len{s'}=1$:] So $s'=\seq{\beta}$ for an action $\beta\in\actions\setminus\set{\tau}$ and there are processes $Q_1',Q_2'\in\procs$ such that $\ec{P\sigma}\bigstep{}\ec{Q_1'}$, $\ec{Q_1'}\transs{\beta}\ec{Q_2'}$, and $\ec{Q_2'}\bigstep{}\ec{Q'}$. \refLem{lem_subst_bigstep_partIII} yields that there is a process $Q_1\in\procs$ with $\ec{P}\bigstep{}\ec{Q_1}$ and $\ec{Q_1\transp{a}{b}}=\ec{Q_1'}$. With $\ec{Q_1\transp{a}{b}}=\ec{Q_1'}\transs{\beta}\ec{Q_2'}$ and \refCor{cor_subst_trans_partIII} respectively \refCor{cor_subst_trans_partIV}, we know that there is an action $\alpha\in\actions$ and a process $Q_2\in\procs$ such that $\ec{Q_1}\transs{\alpha}\ec{Q_2}$ with $\transp{a}{b}(\alpha)=\beta$ and $\ec{Q_2\transp{a}{b}}=\ec{Q_2'}$. \refLem{lem_subst_bigstep_partIII} with $\ec{Q_2\transp{a}{b}}=\ec{Q_2'}\bigstep{}\ec{Q'}$ yields that there is a process $Q\in\procs$ with $\ec{Q_2}\bigstep{}\ec{Q}$ and $\ec{Q\transp{a}{b}}=\ec{Q'}$. Hence, $\ec{P}\bigstep{\seq{\alpha}}\ec{Q}$. Let $s\define\seq{\alpha}$, then also $s\transp{a}{b}=s'$ holds.

\item[Induction hypothesis:] For a number $n\in\N$ and all traces $t'\in\tr$, with $n\leq\len{t'}$ and $a,b\in\names$, $P\in\procs$, $\sigma=\subs{a}{b}$ a substitution with $a\nin\fn{P}$ or $\sigma=\transp{a}{b}$ transposition with $\n{\sigma}\cap\bn{t'}=\emptyset$ and $\ec{P\sigma}\bigstep{t'}\ec{Q'}$, exists a process $Q\in\procs$ and a trace $t\in\tr$ with $\ec{P}\bigstep{t}\ec{Q}$ and $t\transp{a}{b}=t'$ and $\ec{Q\transp{a}{b}}=\ec{Q'}$.

\item[Induction step:] So $\len{s'}=n+1$ and so there is a trace $t'\in\tr$, a process $Q_1'\in\procs$ and an action $\beta\in\actions\setminus\set{\tau}$ with $\ec{P\sigma}\bigstep{t'}\ec{Q_1'}$ and $\ec{Q_1'}\bigstep{\seq{\beta}}\ec{Q'}$ and $s'=\seqconc{t'}{\seq{\beta}}$. The induction hypothesis yields that there is a process $Q_1\in\procs$ and a trace $t\in\tr$ with $\ec{P}\bigstep{t}\ec{Q_1}$ and $t\transp{a}{b}=t'$ and $\ec{Q_1\transp{a}{b}}=\ec{Q_1'}$. Since $\ec{Q_1\transp{a}{b}}=\ec{Q_1'}\bigstep{\seq{\beta}}\ec{Q'}$ we know from the second base case that there is a process $Q\in\procs$ and an action $\alpha\in\actions$ with $\ec{Q_1}\bigstep{\seq{\alpha}}\ec{Q}$ and $\seq{\alpha}\transp{a}{b}=\seq{\beta}$ and $\ec{Q\transp{a}{b}}=\ec{Q'}$. Thus, define $s\define{}\seqconc{t}{\seq{\alpha}}$ then we know $s\transp{a}{b}=s'$ and $\ec{P}\bigstep{s}\ec{Q}$.
\end{description}
Thus, \refLem{lem_subst_bigstep_partIV} is proved by induction over the length of the trace.
\end{prf}

Thereby, we see that the example given after \refLem{lem_subst_bigstep_partII} for the non-existence of the converse is fixed by the previous lemma by the usage of a transposition. Again, consider $P\define\inp{y}{a}$ and $\sigma\define\subs{b}{x}$, then $\ec{P\sigma}=\ec{P}\bigstep{\seq{\inpa{y}{x}}}\ec{\proczero}$. This time, with $s\define\seq{\inpa{y}{b}}$ and $Q\define\proczero$, there is a trace $s\in\tr$ and a process $Q$ with $\ec{P}\bigstep{s}\ec{Q}$ and $s\transp{b}{x}=\seq{\inpa{y}{x}}$ and $\ec{Q\transp{b}{x}}=\ec{\proczero}$.

Since the application of a substitution is only defined on traces with bound names different to the names of the substitution, it is useful to take a closer look at the bound names of a trace. Obviously, we would like to apply a given substitution to an arbitrary trace without destroying anything of its semantics. Intuitively, since those names are bound, we should be able to change the bound names within a trace starting in a process $P$ and the resulting trace is still a valid behavior of $P$.

Furthermore, we would like to abstract from traces where the bound names within are not unique for reasons of simplicity.

\begin{definition}[Unique traces]
\label{def_unique_traces}
We call a trace $t\in\tr$ \findex[uniqueness of bound names!trace]{unique} if 
\begin{enumerate}
\item[(1)] $t_i=\bouta{a}{b}$ implies $\nexists{}c\in\names,j\in\N: i\neq{}j \wedge t_j=\bouta{c}{b}$,
\item[(2)] $t_i=\bouta{a}{b}$ implies $b\nin\n{\seq{t_1,\ldots, t_{i-1}}}$
\end{enumerate}
holds.
\end{definition}

This definition is helpful for determine the scope of a name within a trace. Consider, for example, $P\define{}\procres{c}{\inp{b}{x}.\out{b}{c}.\inp{b}{y}}$. We know $\ec{P}\intrans{b}{a}\ec{\procres{c}{\out{b}{c}.\inp{b}{y}}}\bouttrans{b}{a}\ec{\inp{b}{y}}\intrans{b}{a}\ec{\proczero}$, since $\procres{a}{\out{b}{a}.\inp{b}{y}}\in\ec{\procres{c}{\out{b}{c}.\inp{b}{y}}}$. Thus, $\ec{P}\bigstep{\seq{\inpa{b}{a},\bouta{b}{a},\inpa{b}{a}}}\ec{\proczero}$ with different names $a$. In particular, for $\seq{\inpa{b}{a},\bouta{b}{a},\inpa{b}{a}}$ we know that on the hand the first $a$ is a name visibly to the environment and on the other hand the second one is a different name $a$ and till this moment invisibly for the environment.

It is for this reason that we define the \findex{bound substitution} function, which suitably replaces the bound names within a trace.

\begin{definition}[Bound substitution]
\label{def_bound_substitution}
As $\bnsubstF{}:\pom{\names^2}\times\tr\rightarrow\tr$ we define the \findex{bound substitution} such that\newline
if $\exists{} i\in\N{},x\in\names\nexists{}j\in\N,y\in\names: i<j\wedge{}t_i=\bout{x}{a}\wedge{}t_j=\bout{y}{a}\wedge\len{t}=n$ then
\[\bnsubst{b}{a}{t}=\seqconc{\seq{t_1,\ldots,t_{i-1}}}{\seqconc{\seq{\bout{x}{b}}}{(\seq{t_{i+1},\ldots,t_n}\subs{b}{a})}},\]
if $\exists{} i,j\in\N{},x,y\in\names: i<j\wedge{}t_i=\bout{x}{a}\wedge{}t_j=\bout{y}{a}\wedge\len{t}=n$ then 
\[\bnsubst{b}{a}{t}=\seqconc{\seq{t_1,\ldots,t_{i-1}}}{\seqconc{\seq{\bout{x}{b}}}{\seqconc{(\seq{t_{i+1},\ldots,t_{j-1}}\subs{b}{a})}{\seq{t_j,\ldots,t_n}}}},\]
otherwise
\[\bnsubst{b}{a}{t}=t\]
holds.
\end{definition}

We now show with the previous definition that if there is a big-step with a trace $t$ starting in a process, then there is still a big-step starting in the process for a trace $t'$ if we suitably changed the bound names within the trace $t$ to obtain $t'$.

\begin{lemma}[Big-step semantics under bound name changes]
\label{lem_bn_changes_traces}
\begin{old} %%%%%%%%%%%%%%%%%%%%%%%%%%%%%%%%%%%%%%% OLD preservation for trace semantics %%%%%%%%%%%%%%%%%%%%%%%%%
For a name $a\in\names$ and process $P\in\procs$
\[t\in\traces[\procres{a}{P}] \text{ implies } \bnsubst{b}{a}{t}\in\traces[\procres{a}{P}]\]
holds for all $b\nin\fn{\procres{a}{P}}\cup\n{t}$.
\end{old} %%%%%%%%%%%%%%%%%%%%%%%%%%%%%%%%%%%%%%% OLD end preservation for trace semantics %%%%%%%%%%%%%%%%%%%%%%%
For a name $a\in\names$, a trace $t\in\tr$ and processes $P,Q\in\procs$
\[\ec{P}\bigstep{t}\ec{Q} \text{ implies } \exists{}Q'\in\procs: \ec{P}\bigstep{\bnsubst{b}{a}{t}}\ec{Q'}\]
holds for all $b\nin\fn{P}\cup\n{t}$.
\end{lemma}
\begin{prf}
Let $b\nin\fn{P}\cup\n{t}$, $a\in\names$, $P,Q\in\procs$, and $t\in\tr$ with $\len{t}=n$ for some $n\in\N$, and $\ec{P}\bigstep{t}\ec{Q}$. If $a\nin\bn{t}$ then $t=\bnsubst{b}{a}{t}$. Hence $\ec{P}\bigstep{\bnsubst{b}{a}{t}}\ec{Q}$. If $a\in\bn{t}$ there must be a number $i\in\N$ and a name $x\in\names$ such that $t=\seqconc{t'}{\seqconc{\seq{\bout{x}{a}}}{t''}}$ with $t'\define{}\seq{t_1,\ldots,t_{i-1}}$ and $t''\define\seq{t_{i+1},\ldots,t_n}$. Without loss of generality, we assume that this index is the first occurrence of a bound output with object $a$ within the trace. So there are processes $Q_1,Q_2\in\procs$ with $\ec{P}\bigstep{t'}\ec{Q_1}\transs{\bout{x}{a}}\ec{Q_2}\bigstep{t''}\ec{Q}$. \refLem{lem_bigstep_fn} yields that $\fn{Q_1}\subseteq\fn{P}\cup\n{t'}$ and so $b\nin\fn{Q_1}$. Hence, with \refLem{lem_bn_trans} and since $\fn{\procres{a}{Q_1}}\subseteq\fn{Q_1}$ we know $\ec{Q_1}\transs{\bout{x}{b}}\ec{Q_2\subs{b}{a}}$. If there is no number $j\in\N$ and name $y\in\names$ with $t_j=\bout{y}{a}$ and $j>i$ we know from \refLem{lem_subst_bigstep_partII} that $\ec{Q_2\subs{b}{a}}\bigstep{t''\subs{b}{a}}\ec{Q\subs{b}{a}}$. Furthermore, since in this case $\bnsubst{b}{a}{t}=\seqconc{t'}{\seqconc{\seq{\bout{x}{b}}}{t''\subs{b}{a}}}$ we know $\ec{P}\bigstep{\bnsubst{b}{a}{t}}\ec{Q\subs{b}{a}}$. Otherwise, there is a number $j\in\N$ and a name $y\in\names$ such that $t_j=\bout{y}{a}$ and $j>i$. With $s\define{}\seq{t_{i+1},\ldots,t_{j-1}}$ and $s'\define\seq{t_j,\ldots,t_n}$ we know there is a process $Q_2'\in\procs$ such that $\ec{P}\bigstep{t'}\ec{Q_1}\transs{\bout{x}{a}}\ec{Q_2}\bigstep{s}\ec{Q_2'}\bigstep{s'}\ec{Q}$. Without loss of generality, we assume that this index is the second occurrence of a bound output with object $a$ within the trace. Thus, with \refLem{lem_subst_bigstep_partII} we know $\ec{Q_2\subs{b}{a}}\bigstep{s\subs{b}{a}}\ec{Q_2'\subs{b}{a}}$ holds. If $a\in\fn{Q_2'}$ we know from \refConv{conv_uni_bn} that $a\nin\bn{Q_2'}$. Furthermore, there is no process $R\in\procs$ such that $R=_\alpha{}Q_2'$ and $a\in\bn{R}$, since with $a\in\fn{Q_2'}$ we know $a\in\fn{R}$ for all $R=_\alpha{}Q_2'$ and so the name $a$ would violate \refConv{conv_uni_bn} because $a\in\fn{R}\cup\bn{R}$. But if there is no process $R\in\ec{Q_2'}$ with $a\in\bn{R}$, there is no possibility for a transition $\transs{\bout{y}{a}}$ to start in $\ec{Q_2'}$. Thus, $a\nin\fn{Q_2'}$ and so $\ec{Q_2'\subs{b}{a}}=\ec{Q_2'}$. Hence, $\ec{P}\bigstep{t'}\ec{Q_1}\transs{\bout{x}{b}}\ec{Q_2\subs{b}{a}}\bigstep{s\subs{b}{a}}\ec{Q_2'\subs{b}{a}}\bigstep{s'}\ec{Q}$ and thus, $\ec{P}\bigstep{\bnsubst{b}{a}{t}}\ec{Q}$.
\end{prf}

Thus, we can change the bound names within a trace with any fresh name without changing its semantics. This justifies the following convention.

\begin{conv}[Uniqueness of bound names in traces]
\label{conv_uni_bn_traces}
For a trace $t\in\tr$ and substitutions $\sigma_1,\ldots,\sigma_n$ for some $n\in\N$ under consideration, we stipulate that $t$ is unique and $\n{\sigma_i}\cap\bn{t}=\emptyset$ for every $i\in{}\set{1,\ldots,n}$.
\end{conv}

Note that with this convention the constraint $\n{\sigma}\cap\bn{t}=\emptyset$ in \refLem{lem_subst_bigstep_partII} and \refLem{lem_subst_bigstep_partIV} is no longer harmful.



\chapter{Trace semantics}
\label{sec_de_sem_trace}
% mainfile: ../../Refinement.tex
In the \findex[trace!semantics]{trace semantics} every single trace -- meaning all possible sequences of external behavior -- of a process is collected. Thereby, the internal steps are omitted. Thus, the trace semantics is a \index{semantics!denotational}\findex{denotational semantics}, which represents the communication a process offers to its environment. In \refChap{sec_big-step_semantics} with the big-step semantics we regarded one possible external behavior of a process. Now with the trace semantics we introduce a denotational semantics, which is a collection of all these possible big-steps of a process. Intuitively, the big-step semantics follows \textit{one} way through the transition system produced from the rules of \refFig{fig_ts_early} and collects the visible actions. Whereas the trace semantics collects \textit{all} of those ways through the transition system. This idea is an adaption of the trace semantics for CSP, for example, described in \cite{roscoe}.

%This semantics differs from Roscoe's only by its definition on \picalc{} processes and the corresponding operational semantics.

\section{Definition}
\label{sec_de_sem_trace_def}
% mainfile: ../../Refinement.tex
With the big-step semantics the ability to take an external view of a process' behavior is given. Between every single action in a given sequence there can be as many internal actions as needed. The trace semantics just collects all those traces a process can possibly perform in one set.

\begin{definition}[Trace semantics]
\label{def_trace_semantics}
	For every process $P\in\procs$ the \findex[trace!semantics]{trace semantics} $\mathcal{T}:\procs_\alpha\rightarrow\pom{\tr}$ with
		\[\traces[P] \define \set[\exists Q\in\procs: \ec{P} \bigstep{t} \ec{Q}]{t\in \tr}\]
	is a set of all possible external behavior of the process $P$.
\end{definition}

In this way, processes can be distinguished by their external behavior. We now introduce some properties of this denotational semantics before defining the refinement, which bases upon this semantics.
\begin{old}[Sven says: ist ein gewagter Vergleich] %%%%%%%%%%%%%%%%%%%%%%%%%%%%%% OLD: gewagter Vergleich %%%%%%%%%%%%%%%%%%%%%%%%%%%%%%%%%%%%%%%%%%
Another understanding of the traces could be in the view of languages. Thus, $\traces[P]$ stands for the language which could be retrieved from the automaton which is constructed with the operational semantics where every state is a final state by omitting the internal transitions.
\end{old} %%%%%%%%%%%%%%%%%%%%%%%%%%%%%% end OLD: gewagter Vergleich %%%%%%%%%%%%%%%%%%%%%%%%%%%%%%%%%%%%%%%%%%


\section{Properties}
\label{sec_de_sem_trace_prop}
% mainfile: ../../Refinement.tex
In this section we investigate the trace semantics and some of its properties for gaining results for the refinement defined in \refSec{sec_de_sem_trace_ref}. Furthermore, these properties help in understanding the trace semantics and its applicability. Especially, we examine its behavior while applying a substitution and its relationship to simulation and bisimulation to classify the new semantics in the established context. The main part takes the investigation, whether the trace semantics is compositional or not, resulting by the fact that already the subclass of recursion-free processes is indeed not compositional for every operator.

As a first result we achieve that every set of traces is at most countably infinite. This is because we know that every trace is produced out of a countably infinite set of symbols. In other words, traces are words over a countably infinite alphabet. More precisely, for every $P\in\procs$ we know $\traces[P]\subseteq\left(\names\cup\conames\cup\mathtt{SYMB}\right)^*$ where $\mathtt{SYMB}$ denotes the set containing the angle brackets, the parenthesis, and the comma. Furthermore, we know that the countably infinite sets are closed under union and Kleene star. So, since $\names$ and $\conames$ are countably infinite, we know that the set $\traces[P]$ is countably infinite. Hence, we can define a function $\mathcal{T_I}:\procs_\alpha\rightarrow\pom{\N\times\tr}$, named \findex{indexed trace set}, such that every trace in the trace set of a equivalence class of a process is uniquely mapped to a natural number.

Furthermore, we know that every trace set $\traces[P]$ for a process $P\in\procs$ is closed under the prefix operator. Therefore, consider $s\in\traces[P]$. We know for every prefix $t\in\pref{\set{s}}$ that there is a trace $u\in\tr$ such that $s=\seqconc{t}{u}$. So there is a process $Q\in\procs$ such that $\ec{P}\bigstep{\seqconc{t}{u}}\ec{Q}$ and so $\ec{P}\bigstep{t}\fatsemi\bigstep{u}\ec{Q}$ holds. Hence, there is a process $Q'\in\procs$ such that $\ec{P}\bigstep{t}\ec{Q'}$ and so $t\in\traces[P]$.

Since the empty sequence is a prefix of every trace, it fits that the empty trace is included in every trace set.

\begin{lemma}[Empty trace]
\label{lem_empty_trace}
For all $P\in\procs$
\[\eseq\in\traces[P]\]
holds.
\end{lemma}
\begin{prf}
Let $P\in\procs$. Since $\left(\tautrans\right)^*$ is reflexive, we know that $\ec{P}\left(\tautrans\right)^*\ec{P}$ holds. With the definition of the big-step semantics (\refDef{def_bigstep_semantics}) we get $\ec{P}\bigstep{\eseq}\ec{P}$. With that and from the definition of the trace semantics (\refDef{def_trace_semantics}) we know that $\eseq\in\traces[P]$ holds, since $P\in\procs$.
%%%%%%%%%%%%%%%%%%%%%%%%%%%%%%%%%%%%%%%%%%%%%% START LONG VERSION %%%%%%%%%%%%%%%%%%%%%%%%%%%%%%%%%%%%%%%%%%%%%%%%%%%%%%%%%%%%%%%%%%%%%%%%%%%%%%
\begin{old}{long version}
	Let $P\in\procs$. With the definition of the trace semantics (\refDef{def_trace_semantics}), we know that
		\[\traces[P] = \set[\exists{}Q\in\procs:\ec{P}\bigstep{t}\ec{Q}]{t\in\tr}\]
	holds. Since $P\in\procs$ we have
	\[\traces[P]=\underbrace{\set[\ec{P}\bigstep{t}\ec{P}]{t\in\traces}}_{Y\define} \cup \set[\exists{}Q\in\procs: Q\nin\ec{P} \wedge\ec{P}\bigstep{t}\ec{Q}]{t\in\tr}.\]
	We know from the definition of the big-step semantics (\refDef{def_bigstep_semantics}) that $\ec{P}\bigstep{\eseq}\ec{P}$ if and only if $\ec{P}\left(\tautrans\right)^*\ec{P}$. Since $\left(\tautrans\right)^*$ is a reflexive and transitive closure, we know that $\ec{P}\left(\tautrans\right)^*\ec{P}$ holds, which means $\ec{P}\bigstep{\eseq}\ec{P}$ holds, and therefore $\eseq{}\in{}Y$ and so $\eseq\in\traces[P]$.
\end{old}
%%%%%%%%%%%%%%%%%%%%%%%%%%%%%%%%%%%%%%%%%%%%%% END LONG VERSION %%%%%%%%%%%%%%%%%%%%%%%%%%%%%%%%%%%%%%%%%%%%%%%%%%%%%%%%%%%%%%%%%%%%%%%%%%%%%%
\end{prf}

Now, we examine the behavior of the trace semantics in the context of substitutions.

\subsection{Substitution}
\label{sec_de_sem_trace_prop_subst}
\input{chapters/trace_semantics/trace_substitution}

\subsection{Compositionality}
\label{sec_de_sem_trace_prop_comp}
\input{chapters/trace_semantics/trace_compositionality}

\subsection{Simulation and bisimulation}
\label{_sec_de_sem_trace_sim_bisim}
\input{chapters/trace_semantics/trace_simulation}


\section{Refinement}
\label{sec_de_sem_trace_ref}
% mainfile: ../../Refinement.tex
The definition of traces leads directly to the \findex[trace!refinement]{trace refinement}. Hence, the situation that a process has the same external behavior as another one and, in addition to that, can possibly have some more behavior, is intuitively achieved by set inclusion. This definition directly corresponds to the definition of trace refinement in \gls{CSP}, for example in \cite{roscoe}.

\begin{definition}[\index{refinement}{Trace refinement}]
\label{def_trace_ref}
	Let $P,Q\in\procs$, then $Q$ is a \findex[trace!refinement]{trace refinement} of $P$ ($P\refi Q$) if the inverse set inclusion holds:
		\[P\refi Q \text{ is equivalent to } \traces[Q]\subseteq\traces[P].\]	
	We also say for $P\refi{}Q$ that $Q$ \findex[refinement]{refines} $P$.

	Two processes are called \findex[trace!equivalent]{trace equivalent}, if and only if they have got the same external behavior: $\traces[P]=\traces[Q]$.
\end{definition}

Thus, $Q$ refines $P$ if and only if $Q$ has less behavior than $P$. In contrast to the simulations, where the processes are collected to compare their behavior, with the trace sets we do not observe the processes, but only save the labels of the possible transitions.

Since the refinement is just a set inclusion, we know that all the results of \refSec{sec_de_sem_trace_prop} for set inclusion also hold for the refinement.

\begin{lemma}[Refinement properties]
\label{lemma_ref_properties}
Given two processes $P,Q\in\procs$ and sums $M_1,M_2,M_3,M_4\in\sums$. Then
\begin{itemize}
\item[(1)] For all substitutions $\sigma=\subs{a}{b}$ with $a\nin\fn{P}\cup\fn{Q}$ or $\sigma=\transp{a}{b}$ transposition,
\[Q \refi P \text{ implies } Q\sigma \refi P\sigma\]
holds,
\item[(2)] $Q \refi P$ is equivalent to $\out{a}{x}.Q \refi \out{a}{x}.P$, for all $a,x\in\names$,
\item[(3)] $M_2  \refi M_1$ and $M_4 \refi M_3$ implies $\procchoice{M_2}{M_4} \refi \procchoice{M_1}{M_3}$,
\item[(4)] For $\simu\subseteq\procs_\alpha\times\procs_\alpha$ weak or strong simulation
\[(\ec{P},\ec{Q})\in{}\simu \text{ implies } Q \refi P\]
holds
\end{itemize}
holds.
\end{lemma}
\begin{prf}
\refCor{cor_subst_trace_inclusion}, \refLem{lem_pres_out_choice}, \refLem{lem_weak_sim_trace_inclusion}, \refCor{cor_strong_sim_trace_inclusion} directly yield the corresponding properties.
\end{prf}

In the next section we present some examples of applications of the trace refinement.


\section{Example}
\label{sec_de_sem_trace_example}
% mainfile: ../../Refinement.tex
\begin{figure}[h!]
\lstinputlisting[language=pseudocode, mathescape]{listings/compo_comput.psc}
\caption{Compositional computation of traces $\mathtt{calc}$.}
\label{fig_comp_comput}
\end{figure}
In this section we briefly state the usage of the trace refinement by reconsidering the buffer of \refSec{sec_pi_op_sem}. Furthermore, we introduce a rudimentary idea of a first implementation of the calculation of the trace sets. For this approach we almost directly convert the compositional definitions of \refSec{sec_de_sem_trace_prop_comp} to pseudocode.

\begin{figure}[H]
\lstinputlisting[language=pseudocode, mathescape]{listings/compo_comput_inp.psc}
\caption{Compositional computation of traces $\mathtt{calc}_{\mathtt{inp}}$.}
\label{fig_comp_comput_inp}
\end{figure}

Thus, in \refFig{fig_comp_comput} we list the main calculation algorithm, which decides by the process' structure, which sub-algorithm should be used. All algorithms take a recursion-free process and an finite subset of names as input parameters. The limitation to recursion-free processes depends on the circumstances that we neither have any fix-point algorithm up to now nor showed that one exists. The second input parameter is invented for the calculation of the traces of an input process and denotes the names which can be send from the environment to the system. First of all the process is transformed in extended standard form for a possible application of the algorithm for the parallel composition. The algorithm for an inaction, a choice, and an output process are already given in this listing. They are all the direct translation of the definitions of \refSec{sec_de_sem_trace_prop_comp}.

In \refFig{fig_comp_comput_inp} the algorithm for calculating the traces of an input process is given. In this algorithm we can understand the need of the restriction to a finite set of names. Since all names which can possibly be send from the environment to the process are collected in $\relNames$, we know that only those names and the names which occur free in the process under consideration can possibly establish new communication. Thus, we only have to calculate traces for the processes, where the subject of the input prefix is substituted by such names. To optimize the number of calculations, we take advantage of \refCor{cor_subst_trace_inclusion}.

\begin{figure}[h!]
\lstinputlisting[language=pseudocode, mathescape]{listings/compo_comput_par.psc}
\caption{Compositional computation of traces $\mathtt{calc}_{\mathtt{par}}$.}
\label{fig_comp_comput_par}
\end{figure} %$\mathtt{max}$ maximal deduction steps to reach $T_1'$ and $T_2'$;

The algorithm for calculating the traces of the parallel composition is listed in \refFig{fig_comp_comput_par}. Mainly, the algorithm uses \refDef{def_idx_trace_sets} to calculated the sets of the inductive parallel composition of the traces of $P_1$ and $P_2$. This is also a straightforward translation of this definition.

\begin{figure}[h]
\lstinputlisting[language=pseudocode, mathescape]{listings/compo_comput_res.psc}
\caption{Compositional computation of traces $\mathtt{calc}_{\mathtt{res}}$.}
\label{fig_comp_comput_res}
\end{figure}

\refFig{fig_comp_comput_res} show the computation of the traces of a restriction. This computation directly yields from \refConj{conj_compositionality_traces_res}.

Utilizing the approach, we can investigate the buffer defined in \refSec{sec_pi_op_sem}. Since we do not have any algorithm for computing processes with recursive calls, we have to unfold the examples. For a clearer arrangement we do not write the processes as a whole, but naming parts of the process separately. Those naming do not have to be understood as a recursive call without parameter. First of all we reconsider the one cell buffer of \cite{milner} unfolded once:

\[ONE\_CELL\define\inp{i}{v_1}.\out{o}{v_1}.\inp{i}{v_2}.\out{o}{v_2}.\]
With attention to \refFig{fig_ts_buffer_one_cell} or by using the algorithm, we directly see that 
\[\traces[ONE\_CELL]=\pref{\set[v_1,v_2\in\names]{\seq{\inpa{i}{v_1},\outa{o}{v_1},\inpa{i}{v_2},\outa{o}{v_2}}}}.\]
Furthermore, we can omit the recursive calls of the second buffer and reach
\[TWO\_CELL\define\procchoice{\inp{i}{v_1}.\out{o}{v_1}.\inp{i}{v_2}.\out{o}{v_2}}{\inp{i}{v_1'}.\inp{i}{v_2'}.\out{o}{v_1'}.\out{o}{v_2'}}.\]
Thus, the set of traces is again easily to compute:
\begin{align*}
\begin{array}{rcl}
 \traces[TWO\_CELL]=&\pref{\set[v_1,v_2\in\names]{\seq{\inpa{i}{v_1},\outa{o}{v_1},\inpa{i}{v_2},\out{o}{v_2}}}} \\
\cup&\pref{\set[v_1,v_2\in\names]{\seq{\inpa{i}{v_1},\inpa{i}{v_2},\outa{o}{v_1},\outa{o}{v_2}}}}.
\end{array}
\end{align*}
Thus, since $\traces[ONE\_CELL]\subseteq{}\traces[TWO\_CELL]$ holds, we know that the $ONE\_CELL$ buffer refines the $TWO\_CELL$ buffer. That is $TWO\_CELL \refi ONE\_CELL$. As expected this does not hold the other way round.

We can now think of our invented $TWO\_CELL$ buffer as an implementation of Milner's $FIFO$ buffer defined in \refEx{ex_two_cell_buffer_milner}. Thus, we would like that our implementation satisfy the specification of a FIFO buffer by Milner. Therefore, we unfold the buffer once:
\[FIFO\define\procres[()]{c}{\procpar{\inp{i}{v_1}.\out{c}{v_1}.\inp{i}{v_2}.\out{c}{v_2}}{\inp{c}{v_1'}.\out{o}{v_1'}.\inp{c}{v_2'}.\out{o}{v_2'}}}.\]
So we see, that the channel $c$ is private and hence, only the traces where a communication over $c$ has been established are within the set of traces. All traces may not so obviously be read from the definition of $FIFO$. So we implemented the algorithm presented at the beginning of this section. Thus, we get:
\begin{align*}
\begin{array}{rcl}
 \traces[FIFO]=&\pref{\set[v_1,v_2\in\names]{\seq{\inpa{i}{v_1},\outa{o}{v_1},\inpa{i}{v_2},\out{o}{v_2}}}} \\
\cup&\pref{\set[v_1,v_2\in\names]{\seq{\inpa{i}{v_1},\inpa{i}{v_2},\outa{o}{v_1},\outa{o}{v_2}}}}.
\end{array}
\end{align*}
Hence, we know that the $TWO\_CELL$ and the $FIFO$ buffer are trace equivalent and therefore for this unfolding the $TWO\_CELL$ buffer satisfy its specification. But if we unfold the two cell buffer of \refEx{ex_two_cell_buffer_mine} once and the FIFO buffer triply so that they are both able to perform four read and write passes, we see that we did not developed a FIFO buffer with the same behavior than Milner's one. So, we unfold our two cell buffer once and so define $TC$ with:
\begin{align*}
	\begin{array}{rcl}
TC&\define& \inp{i}{v_1}.\inp{i}{v_2}.\out{o}{v_1}.\out{o}{v_2}.\left(\procchoice{B_1}{B_2}\right) \\
			&+& \inp{i}{v_1'}.\out{o}{v_1'}.\inp{i}{v_2'}.\out{o}{v_2'}.\left(\procchoice{B_3}{B_4}\right) \\
B_1&\define&\inp{i}{v_3}.\inp{i}{v_4}.\out{o}{v_3}.\out{o}{v_4} \\
B_2&\define&\inp{i}{v_3'}.\out{o}{v_3'}.\inp{i}{v_4'}.\out{o}{v_4'} \\
B_3&\define&\inp{i}{v_3''}.\inp{i}{v_4''}.\out{o}{v_3''}.\out{o}{v_4''} \\
B_4&\define&\inp{i}{v_3'''}.\out{o}{v_3'''}.\inp{i}{v_4'''}.\out{o}{v_4'''}.
\end{array}
\end{align*}
Furthermore, with $F_2$ we denote the triply unfolded FIFO buffer with:
\begin{align*}
\begin{array}{rcl}
F_2&\define&\procres[()]{c}{\procpar{\inp{i}{v_1}.\out{c}{v_1}.\inp{i}{v_2}.\out{c}{v_2}.B_1}{\inp{c}{v_1'}.\out{o}{v_1'}.\inp{c}{v_2'}.\out{o}{v_2'}.B_2}}\\
B_1&\define&\inp{i}{v_3}.\out{c}{v_3}.\inp{i}{v_4}.\out{c}{v_4} \\
B_2&\define&\inp{c}{v_3'}.\out{o}{v_3'}.\inp{c}{v_4'}.\out{o}{v_4'}.
\end{array}
\end{align*}
The traces for the two cell buffer $TC$ are again readable of the definition, since the traces of the choice operator are just the union of the separate parts. Thus, we know
\begin{align*}
	\begin{array}{rcl}
		\traces[TC]&=&\pref{\set[v_1,\ldots{},v_4\in\names]{\seq{\inpa{i}{v_1},\inpa{i}{v_2},\outa{o}{v_1},\outa{o}{v_2},\inpa{i}{v_3},\inpa{i}{v_4},\outa{o}{v_3},\outa{o}{v_4}}}} \\
				&\cup&\pref{\set[v_1,\ldots{},v_4\in\names]{\seq{\inpa{i}{v_1},\inpa{i}{v_2},\outa{o}{v_1},\outa{o}{v_2},\inpa{i}{v_3},\outa{o}{v_3},\inpa{i}{v_4},\outa{o}{v_4}}}} \\
				&\cup&\pref{\set[v_1,\ldots{},v_4\in\names]{\seq{\inpa{i}{v_1},\outa{o}{v_1},\inpa{i}{v_2},\outa{o}{v_2},\inpa{i}{v_3},\inpa{i}{v_4},\outa{o}{v_3},\outa{o}{v_4}}}} \\
				&\cup&\pref{\set[v_1,\ldots{},v_4\in\names]{\seq{\inpa{i}{v_1},\outa{o}{v_1},\inpa{i}{v_2},\outa{o}{v_2},\inpa{i}{v_3},\outa{o}{v_3},\inpa{i}{v_4},\outa{o}{v_4}}}} 
	\end{array}
\end{align*}
holds. For the FIFO buffer $F_2$ we needed to adapt our rudimentary implementation a bit for this special case. Otherwise the computer needed more than $35$GB RAM (including swap), although we swap all actual not needed sets of the inductive parallel composition to the hard drive. Thus, since we know that the input variables never are used as channels we always take the first case in \refFig{fig_comp_comput_inp}. That is, we are treating the input variables in any case as not free within the process. Furthermore, we just compute the traces for four names possible as input variable. At least, for the parallel composition we added another condition for adding traces to the sets of the inductive parallel composition. Since we know that after computing the traces of the parallel composition, we will in this case filter out all traces containing the name $c$ with the restriction algorithm. Thus, we already omitting them during the computation of the parallel composition. This approach already calculates $463761$ traces for the inductive parallel composition with index $0$ and ends up by $24505$ traces. This all leads to the following result of the set of traces for the $F_2$ FIFO buffer:
\begin{align*}
	\begin{array}{rcl}
		\traces[F_2]&=&\pref{\set[v_1,\ldots{},v_4\in\names]{\seq{\inpa{i}{v_1},\outa{o}{v_1},\inpa{i}{v_2},\outa{o}{v_2},\inpa{i}{v_3},\outa{o}{v_3},\inpa{i}{v_4},\outa{o}{v_4}}}} \\
&\cup&\pref{\set[v_1,\ldots{},v_4\in\names]{\seq{\inpa{i}{v_1},\inpa{i}{v_2},\outa{o}{v_1},\inpa{i}{v_3},\outa{o}{v_2},\inpa{i}{v_4},\outa{o}{v_3},\outa{o}{v_4}}}} \\
&\cup&\pref{\set[v_1,\ldots{},v_4\in\names]{\seq{\inpa{i}{v_1},\inpa{i}{v_2},\outa{o}{v_1},\inpa{i}{v_3},\outa{o}{v_2},\outa{o}{v_3},\inpa{i}{v_4},\outa{o}{v_4}}}} \\
&\cup&\pref{\set[v_1,\ldots{},v_4\in\names]{\seq{\inpa{i}{v_1},\outa{o}{v_1},\inpa{i}{v_2},\outa{o}{v_2},\inpa{i}{v_3},\inpa{i}{v_4},\outa{o}{v_3},\outa{o}{v_4}}}} \\
&\cup&\pref{\set[v_1,\ldots{},v_4\in\names]{\seq{\inpa{i}{v_1},\inpa{i}{v_2},\outa{o}{v_1},\outa{o}{v_2},\inpa{i}{v_3},\inpa{i}{v_4},\outa{o}{v_3},\outa{o}{v_4}}}} \\
&\cup&\pref{\set[v_1,\ldots{},v_4\in\names]{\seq{\inpa{i}{v_1},\outa{o}{v_1},\inpa{i}{v_2},\inpa{i}{v_3},\outa{o}{v_2},\inpa{i}{v_4},\outa{o}{v_3},\outa{o}{v_4}}}} \\
&\cup&\pref{\set[v_1,\ldots{},v_4\in\names]{\seq{\inpa{i}{v_1},\inpa{i}{v_2},\outa{o}{v_1},\outa{o}{v_2},\inpa{i}{v_3},\outa{o}{v_3},\inpa{i}{v_4},\outa{o}{v_4}}}} \\
&\cup&\pref{\set[v_1,\ldots{},v_4\in\names]{\seq{\inpa{i}{v_1},\outa{o}{v_1},\inpa{i}{v_2},\inpa{i}{v_3},\outa{o}{v_2},\outa{o}{v_3},\inpa{i}{v_4},\outa{o}{v_4}}}}. 
	\end{array}
\end{align*}

Thus, we get that $F_2$ does satisfy the specification of a FIFO buffer, but since for example $\seq{\inpa{i}{v_1},\inpa{i}{v_2},\outa{o}{v_1},\inpa{i}{v_3},\outa{o}{v_2},\outa{o}{v_3},\inpa{i}{v_4},\outa{o}{v_4}}\nin\traces[TC]$, we know those buffer are not trace equivalent.

For a buffer not satisfying the specification of Milner's FIFO buffer, we consider
\begin{align*}
\begin{array}{rcl}
	BUF&\define&\inp{i}{v_1}.\out{o}{v_1}.\inp{i}{v_2}.\out{o}{v_2}\\
			& +& \inp{i}{v_1'}.\inp{i}{v_2'}.\out{o}{v_1'}.\out{o}{v_2'}\\
			& +& \inp{i}{v_1''}.\inp{i}{v_2''}.\out{o}{v_2''}.\out{o}{v_1''}.
\end{array}
\end{align*}
Thus, we know
\begin{align*}
\begin{array}{rcl}
 \traces[BUF]=&\pref{\set[v_1,v_2\in\names]{\seq{\inpa{i}{v_1},\outa{o}{v_1},\inpa{i}{v_2},\out{o}{v_2}}}} \\
\cup&\pref{\set[v_1,v_2\in\names]{\seq{\inpa{i}{v_1},\inpa{i}{v_2},\outa{o}{v_1},\outa{o}{v_2}}}}\\
\cup&\pref{\set[v_1,v_2\in\names]{\seq{\inpa{i}{v_1},\inpa{i}{v_2},\outa{o}{v_2},\outa{o}{v_1}}}}.
\end{array}
\end{align*}
holds. Hence, $BUF$ is not satisfying the specification of $FIFO$, since for example $\seq{\inpa{i}{v_1},\inpa{i}{v_2},\outa{o}{v_2},\outa{o}{v_1}}\nin\traces[FIFO]$.


%%%%%%%%%%%%%%%%%%%%%%%%%%%%%%%%%%%%%%%%%%%%%%%%%%%%%%%%%% START SOME OLD EXAMPLES %%%%%%%%%%%%%%%%%%%%%%%%%%%%%%%%%%%%%%%%%%%%%%%%%%%%%%%%%%%%%%%%



\begin{old}{example of operational semantics chapter with things needed for refinement}
%  * <i valueB,i valueA,o_<valueB>,o_<valueA>>,
%         * <i valueA,i valueB,o_<valueA>,o_<valueB>>,
%         * <i valueA,o_<valueA>,i valueB,o_<valueB>>,
%         * <i valueB,o_<valueB>,i valueA,o_<valueA>>
\begin{align*}
  \begin{array}{rcl}
     B   &\define{}&\procres[()]{c_1}{\procpar{B_1}{B_2}} \\
     B_1 &\define{}&\inp{i}{v_1}.\out{c_1}{v_1}.B_3 \\
     B_2 &\define{}&\inp{c_1}{v_2}.\out{o}{v_2}.B_4 
  \end{array}
\end{align*}

\begin{equation*}	
	\begin{aligned}	
  \begin{array}{rcl}
     B_3 &\define{}&\procres[()]{c_2}{\procpar{B_5}{B_6}} \\
     B_5 &\define{}&\inp{i}{v_3}.\out{c_2}{v_3} \\
     B_6 &\define{}&\inp{c_2}{v_4}.\out{c_1}{v_4} 
  \end{array}
\end{aligned}\quad\quad
\begin{aligned}
  \begin{array}{rcl}
     B_4   &\define{}&\procres[()]{c_3}{\procpar{B_7}{B_8}} \\
     B_7 &\define{}&\inp{c_1}{v_5}.\out{c_3}{v_5} \\
     B_8 &\define{}&\inp{c_3}{v_6}.\out{o}{v_6} 
  \end{array}
\end{aligned}
\end{equation*}
Thus, 
\begin{align*}
B\struc\procres{}{}\;c_1,c_2,c_3\bigl(\inp{i}{v_1}.\out{c_1}{v_1}.\bigl(\procpar{\inp{i}{v_3}.\out{c_2}{v_3} }{\inp{c_2}{v_4}.\out{c_1}{v_4}}\bigr) \\
\;\mid\; \inp{c_1}{v_2}.\out{o}{v_2}.\bigl(\procpar{\inp{c_1}{v_5}.\out{c_3}{v_5}}{\inp{c_3}{v_6}.\out{o}{v_6}}\bigr)\bigr)
\end{align*}

thewn@apu:~/masterarbeit/git/impl/REF/dist java -Xmx50g -jar REF.jar 
P0\define{}\procres[()]{c1}{P0'\define{}\procpar{P1\define{}\inp{i}{v1}.P1'\define{}\out{c1}{v1}.P3\define{}\procres[()]{c2}{P3'\define{}\procpar{P5\define{}\inp{i}{v3}.P5'\define{}\out{c2}{v3}.P9\define{}\procres[()]{c2}{P9'\define{}\procpar{P13\define{}\inp{i}{v7}.P13'\define{}\out{c4}{v7}.\proczero{}}{P14\define{}\inp{c4}{v8}.P14'\define{}\out{c2}{v8}.\proczero{}}}}{P6\define{}\inp{c2}{v4}.P6'\define{}\out{c1}{v4}.P10\define{}\procres[()]{c3}{P10'\define{}\procpar{P15\define{}\inp{c2}{v9}.P15'\define{}\out{c5}{v9}.\proczero{}}{P16\define{}\inp{c5}{v10}.P16'\define{}\out{c1}{v10}.\proczero{}}}}}}{P2\define{}\inp{c1}{v2}.P2'\define{}\out{o}{v2}.P4\define{}\procres[()]{c3}{P4'\define{}\procpar{P7\define{}\inp{c1}{v5}.P7'\define{}\out{c3}{v5}.P11\define{}\procres[()]{c2}{P11'\define{}\procpar{P17\define{}\inp{c1}{v11}.P17'\define{}\out{c6}{v11}.\proczero{}}{P18\define{}\inp{c6}{v12}.P18'\define{}\out{c3}{v12}.\proczero{}}}}{P8\define{}\inp{c3}{v6}.P8'\define{}\out{o}{v6}.P12\define{}\procres[()]{c3}{P12'\define{}\procpar{P19\define{}\inp{c3}{v13}.P19'\define{}\out{c7}{v13}.\proczero{}}{P20\define{}\inp{c7}{v14}.20'\define{}\out{o}{v14}.\proczero{}}}}}}}
P1\define{}\inp{i}{v1}.P1'\define{}\out{c1}{v1}.P3\define{}\procres[()]{c2}{P3'\define{}\procpar{P5\define{}\inp{i}{v3}.P5'\define{}\out{c2}{v3}.P9\define{}\procres[()]{c2}{P9'\define{}\procpar{P13\define{}\inp{i}{v7}.P13'\define{}\out{c4}{v7}.\proczero{}}{P14\define{}\inp{c4}{v8}.P14'\define{}\out{c2}{v8}.\proczero{}}}}{P6\define{}\inp{c2}{v4}.P6'\define{}\out{c1}{v4}.P10\define{}\procres[()]{c3}{P10'\define{}\procpar{P15\define{}\inp{c2}{v9}.P15'\define{}\out{c5}{v9}.\proczero{}}{P16\define{}\inp{c5}{v10}.P16'\define{}\out{c1}{v10}.\proczero{}}}}}
P2\define{}\inp{c1}{v2}.P2'\define{}\out{o}{v2}.P4\define{}\procres[()]{c3}{P4'\define{}\procpar{P7\define{}\inp{c1}{v5}.P7'\define{}\out{c3}{v5}.P11\define{}\procres[()]{c2}{P11'\define{}\procpar{P17\define{}\inp{c1}{v11}.P17'\define{}\out{c6}{v11}.\proczero{}}{P18\define{}\inp{c6}{v12}.P18'\define{}\out{c3}{v12}.\proczero{}}}}{P8\define{}\inp{c3}{v6}.P8'\define{}\out{o}{v6}.P12\define{}\procres[()]{c3}{P12'\define{}\procpar{P19\define{}\inp{c3}{v13}.P19'\define{}\out{c7}{v13}.\proczero{}}{P20\define{}\inp{c7}{v14}.20'\define{}\out{o}{v14}.\proczero{}}}}}
P3\define{}\procres[()]{c2}{P3'\define{}\procpar{P5\define{}\inp{i}{v3}.P5'\define{}\out{c2}{v3}.P9\define{}\procres[()]{c2}{P9'\define{}\procpar{P13\define{}\inp{i}{v7}.P13'\define{}\out{c4}{v7}.\proczero{}}{P14\define{}\inp{c4}{v8}.P14'\define{}\out{c2}{v8}.\proczero{}}}}{P6\define{}\inp{c2}{v4}.P6'\define{}\out{c1}{v4}.P10\define{}\procres[()]{c3}{P10'\define{}\procpar{P15\define{}\inp{c2}{v9}.P15'\define{}\out{c5}{v9}.\proczero{}}{P16\define{}\inp{c5}{v10}.P16'\define{}\out{c1}{v10}.\proczero{}}}}}
P4\define{}\procres[()]{c3}{P4'\define{}\procpar{P7\define{}\inp{c1}{v5}.P7'\define{}\out{c3}{v5}.P11\define{}\procres[()]{c2}{P11'\define{}\procpar{P17\define{}\inp{c1}{v11}.P17'\define{}\out{c6}{v11}.\proczero{}}{P18\define{}\inp{c6}{v12}.P18'\define{}\out{c3}{v12}.\proczero{}}}}{P8\define{}\inp{c3}{v6}.P8'\define{}\out{o}{v6}.P12\define{}\procres[()]{c3}{P12'\define{}\procpar{P19\define{}\inp{c3}{v13}.P19'\define{}\out{c7}{v13}.\proczero{}}{P20\define{}\inp{c7}{v14}.20'\define{}\out{o}{v14}.\proczero{}}}}}
P5\define{}\inp{i}{v3}.P5'\define{}\out{c2}{v3}.P9\define{}\procres[()]{c2}{P9'\define{}\procpar{P13\define{}\inp{i}{v7}.P13'\define{}\out{c4}{v7}.\proczero{}}{P14\define{}\inp{c4}{v8}.P14'\define{}\out{c2}{v8}.\proczero{}}}
P6\define{}\inp{c2}{v4}.P6'\define{}\out{c1}{v4}.P10\define{}\procres[()]{c3}{P10'\define{}\procpar{P15\define{}\inp{c2}{v9}.P15'\define{}\out{c5}{v9}.\proczero{}}{P16\define{}\inp{c5}{v10}.P16'\define{}\out{c1}{v10}.\proczero{}}}
P7\define{}\inp{c1}{v5}.P7'\define{}\out{c3}{v5}.P11\define{}\procres[()]{c2}{P11'\define{}\procpar{P17\define{}\inp{c1}{v11}.P17'\define{}\out{c6}{v11}.\proczero{}}{P18\define{}\inp{c6}{v12}.P18'\define{}\out{c3}{v12}.\proczero{}}}
P8\define{}\inp{c3}{v6}.P8'\define{}\out{o}{v6}.P12\define{}\procres[()]{c3}{P12'\define{}\procpar{P19\define{}\inp{c3}{v13}.P19'\define{}\out{c7}{v13}.\proczero{}}{P20\define{}\inp{c7}{v14}.20'\define{}\out{o}{v14}.\proczero{}}}
Loop startet: 1times.
Calculation idx_0 for 1 and 1
idx_0: 1
idx_i: 2
idx: 2
idx_i: 5
idx: 7
idx_i: 7
idx: 14
idx_i: 6
idx: 20
idx_i: 0
idx: 20
Final trace count: 20
Restriction startet: 1times.
Restriction of 20 traces.
Resultin in: 9 traces.
Loop startet: 2times.
Calculation idx_0 for 1 and 1
idx_0: 1
idx_i: 2
idx: 2
idx_i: 5
idx: 7
idx_i: 7
idx: 14
idx_i: 6
idx: 20
idx_i: 0
idx: 20
Final trace count: 20
Restriction startet: 2times.
Restriction of 20 traces.
Resultin in: 20 traces.
Loop startet: 3times.
Calculation idx_0 for 9 and 20
idx_0: 180
idx_i: 360
idx: 360
idx_i: 900
idx: 1260
idx_i: 1800
idx: 3060
idx_i: 3750
idx: 6810
idx_i: 6966
idx: 13776
idx_i: 11697
idx: 25473
idx_i: 17442
idx: 42915
idx_i: 22494
idx: 65409
idx_i: 23730
idx: 89139
idx_i: 18480
idx: 107619
idx_i: 8316
idx: 115935
idx_i: 0
idx: 115935
Final trace count: 115935
Restriction startet: 3times.
Restriction of 115935 traces.
Resultin in: 3273 traces.
Loop startet: 4times.
Calculation idx_0 for 1 and 1
idx_0: 1
idx_i: 2
idx: 2
idx_i: 5
idx: 7
idx_i: 7
idx: 14
idx_i: 6
idx: 20
idx_i: 0
idx: 20
Final trace count: 20
Restriction startet: 4times.
Restriction of 20 traces.
Resultin in: 20 traces.
Loop startet: 5times.
Calculation idx_0 for 1 and 1
idx_0: 1
idx_i: 2
idx: 2
idx_i: 5
idx: 7
idx_i: 7
idx: 14
idx_i: 6
idx: 20
idx_i: 0
idx: 20
Final trace count: 20
Restriction startet: 5times.
Restriction of 20 traces.
Resultin in: 2 traces.
Loop startet: 6times.
Calculation idx_0 for 20 and 2
idx_0: 40
idx_i: 80
idx: 80
idx_i: 200
idx: 280
idx_i: 400
idx: 680
idx_i: 776
idx: 1456
idx_i: 1380
idx: 2836
idx_i: 2194
idx: 5030
idx_i: 2992
idx: 8022
idx_i: 3332
idx: 11354
idx_i: 2751
idx: 14105
idx_i: 1260
idx: 15365
idx_i: 0
idx: 15365
Final trace count: 15365
Restriction startet: 6times.
Restriction of 15365 traces.
Resultin in: 1190 traces.
Loop startet: 7times.
Calculation idx_0 for 3273 and 1190
idx_0: 3894870
idx_i: 7789740
idx: 7789740
idx_i: 19474350
idx: 27264090
idx_i: 38948700
idx: 66212790
idx_i: 90080580
Killed


\input{./images/transitionsystems/op_sem_buffer_one_cell}
To explain this definition, it is helpful to reactivate the definitions of a buffer with one and with two cells of \refChap{sec_preliminaries} and their operational semantics of \refChap{sec_pi_op_sem}. 

First of all, the one cell buffer $B_1$, which was defined by Milner in \cite{milner}, with one value $val$ and two channels $in$ and $out$ to store respectively loads the value:
\begin{align*}
	\procdef{B_1}{in,out}&\define\inp{in}{val}.\proccall{O_1}{val,in,out} \\
	\procdef{O_1}{val,in,out}&\define\out{out}{val}.\proccall{B_1}{in,out}
\end{align*}
Since the buffer uses recursive calls to be ready again after one pass of storage and loading, the traces have to be infinite. With attention to \refFig{fig_ts_buffer_one_cell} for the traces of the one cell buffer holds:
\[\traces[\procdef{B_1}{in,out}] = \pref{\set[n\in\names]{\seq{\inpa{in}{n},\outa{out}{n}}}^*}\]
If the buffer is extended to a buffer with two cells and the fifo strategy, one possible modeling could be process $B_2$:
\begin{align*}
	\procdef{B_2}{in,out}&\define{\procchoice{\proccall{C_1}{in,out}}{\proccall{C_2}{in,out}}}  \\
	\procdef{C_1}{in,out}&\define\inp{in}{val_1}.\inp{in}{val_2}.\out{out}{val_1}.\out{out}{val_2}.\proccall{B_2}{in,out}	\\
	\procdef{C_2}{in,out}&\define\inp{in}{val_1}.\out{out}{val_1}.\proccall{B_2}{in,out}	
\end{align*}
\input{./images/transitionsystems/op_sem_buffer_fifo_1}
Since this buffer can also use just one cell, it has the same behavior as $B_1$; but with its second cell there is a bit more behavior added. With the properties of the trace semantics or a look at \refFig{fig_ts_buffer_fifo_1} the trace semantics of buffer $B_2$ evolves to:
\begin{align*}
	\traces[\procdef{B_2}{in,out}]=&\pref{\set[n\in\names]{\seq{\inpa{in}{n},\outa{out}{n}}}^*} \\
						\cup &\pref{\set[x,y\in\names]{\seq{\inpa{in}{x},\inpa{in}{y},\outa{out}{x},\outa{out}{y}}}^*}.
\end{align*}
So $B_1$ refines $B_2$ ($\procdef{B_2}{in,out}\refi\procdef{B_1}{in,out}$), because $\traces[\procdef{B_1}{in,out}] \subset \traces[\procdef{B_2}{in,out}]$.

It is also possible to define a fifo buffer with two cells, as a linked list:
\begin{align*}
	\procdef{FIFO}{in,out}&\define\procres[a]{com}{\procpar{\proccall{B_3}{in,com}}{\proccall{B_3}{com,out}}}\\
	\procdef{B_3}{in,out}&\define\inp{in}{val}.\proccall{O_2}{val,in,out} \\
	\procdef{O_2}{val,in,out}&\define\out{out}{val}.\proccall{B_3}{in,out}
\end{align*}
\todo{this should not be true, just $FIFO\refi{}B_2$.}
Those two fifo buffer are trace equivalent ($\traces[\procdef{B_2}{in,out}]=\traces[\procdef{FIFO}{in,out}]$). This could be helpful if one buffer has a proven property over its external behavior, then it is also known, that the other process fulfills the same property.
\todo[inline]{for finite traces adapted milners example with a parallel process modelling a counter und adding an output value for the recursive call}
\end{old}
%%%%%%%%%%%%%%%%%%%%%%%%%%%%%%%%%%%%%%%%%%%%%%%%%%%%%%%%%%% END SOME OLD EXAMPLES %%%%%%%%%%%%%%%%%%%%%%%%%%%%%%%%%%%%%%%%%%%%%%%%%%%%%%%%%%%%%%%%



%\chapter{Denotational semantics}
%\label{sec_de_semantics}
%% mainfile: ../Refinement.tex
In this chapter we introduce three denotational semantics for \picalc{} processes. The definitions are very similar to the denotational semantics of \gls{CSP} processes defined by Bill Roscoe in \cite{roscoe}. With attention to an easy definition of the refinement of \picalc{} processes, the denotational semantics consists of sets retrieved from the operational semantics presented in \refSec{sec_pi_op_sem}.

Therefore, in \refSec{sec_de_sem_trace} we define a trace semantics $\traces$, which describes a process' behavior regarding its external, neglecting its internal ability. The internal behavior is taken to account in the failures semantics $\failures$ presented in \refSec{sec_de_sem_failures} and a somewhat different view on internal divergence of a process is constituted with the failures-divergence semantics $\fd$ in \refSec{sec_de_sem_fd}.

\section{Trace semantics}
\label{sec_de_sem_trace}
% mainfile: ../../Refinement.tex
In the \findex[trace!semantics]{trace semantics} every single trace -- meaning all possible sequences of external behavior -- of a process is collected. Thereby, the internal steps are omitted. Thus, the trace semantics is a \index{semantics!denotational}\findex{denotational semantics}, which represents the communication a process offers to its environment. In \refChap{sec_big-step_semantics} with the big-step semantics we regarded one possible external behavior of a process. Now with the trace semantics we introduce a denotational semantics, which is a collection of all these possible big-steps of a process. Intuitively, the big-step semantics follows \textit{one} way through the transition system produced from the rules of \refFig{fig_ts_early} and collects the visible actions. Whereas the trace semantics collects \textit{all} of those ways through the transition system. This idea is an adaption of the trace semantics for CSP, for example, described in \cite{roscoe}.

%This semantics differs from Roscoe's only by its definition on \picalc{} processes and the corresponding operational semantics.

\section{Definition}
\label{sec_de_sem_trace_def}
\input{./chapters/trace_semantics/trace_definition}

\section{Properties}
\label{sec_de_sem_trace_prop}
\input{./chapters/trace_semantics/trace_properties}

\section{Refinement}
\label{sec_de_sem_trace_ref}
\input{./chapters/trace_semantics/trace_refinement}

\section{Example}
\label{sec_de_sem_trace_example}
\input{./chapters/trace_semantics/trace_examples}


%\section{Failures semantics}
%\label{sec_de_sem_failures}
%% mainfile: ../Refinement.tex
\todo{instead of just traces is not nice}The failures semantics collects instead of just traces, tuple of traces and a set of so-called failures. These failure sets describe for the process reachable with the respected trace and if no internal transition is possible, which transitions cannot be used for the next step. This takes the internal behavior of a process into account.

Like the trace semantics, the failures semantics do not really differ from the definition of Roscoe and this section has a similar structure as the previous one.
\subsection{Definition}
\label{sec_de_sem_failures_def}

\subsection{Examples}
\label{sec_de_sem_failures_bsp}
\todo[inline]{example why not just adding taus to the traces}

\subsection{Properties}
\label{sec_de_sem_failures_prop}

\subsection{Refinement}
\label{sec_de_sem_failures_ref}



%\section{Failures-divergence semantics}
%\label{sec_de_sem_fd}
%% mainfile: ../Refinement.tex
\subsection{Definition}
\label{sec_de_sem_fd_def}

\subsection{Properties}
\label{sec_de_sem_fd_prop}

\subsection{Refinement}
\label{sec_de_sem_fd_ref}


%\section{Related work}
%\label{sec_de_sem_rel_work}
%\todo{Relation for inheritance of Canal et al. adaption to refinement?}


%\chapter{Reduction of semantics}
%\label{sec_reduc_semantics}
%% mainfile: ../Refinement.tex
By reason of the dimension of the transition systems resulting already from small processes, \refSec{sec_op_sem_coverage} introduces an adaption of Sangiorgi's semantics and shows the equivalence of both. This notion with the associated adaption of the denotational semantics also helps to simplify the proofs, that one process is a refinement of an other.
\section{Operational semantics}
\label{sec_op_sem_coverage}
For this section we follow this notion but for the sense of simplification for the denotational semantics we get rid of them in \refSec{sec_op_sem_coverage} for the rest of the thesis.\todo{alpha conversion equiv classes of processes}
\subsection{Definition}
\label{sec_op_sem_cov_def}
% mainfile: ../../Refinement.tex
\begin{figure}
\begin{gather*}
\kalRule{TAU}{}{}{\tau.P \trans{\tau} P} \\\\
\kalRule{OUT}{}{}{\out{x}{y}.P \trans{\out{x}{y}} P} \quad\quad \kalRule{IN}{}{}{\inp{x}{z}.P \transin{x}{z} P} \\\\
\kalRule{SUM_L}{}{P \trans{\alpha} P'}{P+Q \trans{\alpha} P'} \quad\quad\kalRule{SUM_R}{}{P \trans{\alpha} P'}{Q+P \trans{\alpha} P'} \\\\
\kalRule{PAR_L}{P \transin{a}{x} P'}{z\in\fn{P}\cup\fn{Q}\cup\set{x}}{\procpar{P}{Q} \transin{a}{z} \procpar{P'\subs{x}{z}}{Q}} \quad\quad \kalRule{PAR_R}{}{P \trans{\alpha} P'}{\procpar{Q}{P} \trans{\alpha} \procpar{Q}{P'}} \\\\
\kalRule[x\nin n(\alpha)]{RES}{}{P \trans{\alpha} P'}{\procres{x}{P} \trans{\alpha} \procres{x}{P'}} \quad\quad \kalRule[z\neq x]{OPEN}{}{P \transout{x}{z} P'}{\procres{z}{P} \transbout{x}{z}P'} \\\\
\kalRule{COM_L}{P \transin{x}{y} P'}{Q \transout{x}{z} Q'}{\procpar{P}{Q} \transtau \procpar{P'\subs{z}{y}}{Q'}} \quad\quad\kalRule{COM_R}{P \transin{x}{y} P'}{Q \transout{x}{z} Q'}{\procpar{Q}{P} \transtau \procpar{Q'}{P'\subs{z}{y}}} \\\\
\kalRule{CLOSE_L}{P \transbout{x}{z} P'}{Q \transin{x}{z} Q'}{\procpar{P}{Q} \transtau \procres[a]{z}{\procpar{P'}{Q'}}} \quad\quad\kalRule{CLOSE_R}{P \transbout{x}{z} P'}{Q \transin{x}{z} Q'}{\procpar{Q}{P} \transtau \procres[a]{z}{\procpar{Q'}{P'}}}\\\\
\kalRule[\procdef{A}{\parl{w}}:=P]{CALL}{}{}{\proccall{A}{\parl{v}} \transtau P\subs{\parl{v}}{\parl{w}}} 
\end{gather*}
\label{fig_ts_mine}
\caption{Transition system for $\pi$-calculus}
\end{figure}


\subsection{Properties and proof}
\label{sec_op_sem_cov_prop}
\begin{theorem}[Coverability]
	Sangiorgi's transition system mentioned in \refFig{fig_ts_early} can be covered by the transition system specified in \refFig{fig_ts_mine}; meaning:
	for all $P,Q\in\procs$ and exists $a,x\in\names$ hold
	\begin{enumerate}
		\item $\ec{P}\tautrans\ec{Q} \gdw P\transtau Q$
		\item $\ec{P}\outtrans{a}{x}\ec{Q} \gdw P\transout{a}{x} Q$
		\item $\ec{P}\intrans{a}{x}\ec{Q} \gdw \exists y\in\fn{P}\cup\set{x}: P\transin{a}{y} Q$
		\item $\ec{P}\bouttrans{a}{x}\ec{Q} \gdw \exists y\in\fn{P}\cup\set{x}: P\transbout{a}{y} Q$
	\end{enumerate}
\end{theorem}

\section{Denotational semantics}
\label{sec_op_sem_coverage}

\begin{align*}
\underbrace{\seq{s_1,\dots,s_n}}_{=_{def}seq_1} \simeq \underbrace{\seq{s_1',\dots,s_m'}}_{=_{def} seq_2} \gdw & \exists l,o,k\in\N\; with\; l+o=k: \\
& \exists x_1,\dots,x_k\in\names\setminus\left(\fn{seq_1}\cup\fn{seq_2}\right) \\
& \exists y_1,\dots,y_l\in\bn{seq_1} \exists z_1,\dots,z_o\in\bn{seq_2}: \\
& seq_1\subs{x_1,\dots,x_l}{y_1,\dots,y_l} = seq_2\subs{x_{l+1},\dots,x_k}{z_1,\dots,z_o}
\end{align*}

\[\traces[P] \subsetsim \traces[Q] \gdw \forall s\in\traces[P]\exists s'\in\traces[Q]: s\simeq s' \]




	% mainfile: ../Refinement.tex
\chapter{Conclusion and future work}
\label{sec_conclusion}
In this thesis ...


	\singlespacing
	%% ------------------------------------------------- Anhang -------------------------------------------------------------------
	%% Anhang
	%\appendix
	%
\section{Tabellen}
\begin{landscape}
	\begin{table}[htb]
	\caption{Credentials der PG-RiO}
	\begin{tabular}{|p{0.15\linewidth}|p{0.20\linewidth}|p{0.60\linewidth}|}
		\hline
		Zweck & Benutzername & Passwort \\
		\hline
		Dockerhub & pgrio & ~\newline \\
		\hline
		VMs & pgrio & ~\newline \\
		\hline
		SK"=Verwaltung & svadmin & ~\newline \\
		\hline
		Jenkins & pgrio & ~\newline \\
		\hline
		DB"=Writer & writer & ~\newline \\
		\hline
		DB"=Admin & admin & ~\newline \\
		\hline
		MQTT"=Datacollector & pgrio"=datacollector & ~\newline \\
		\hline
		API"=Routing & routing & ~\newline \\
		\hline
		App"=Signierung Google Playstore & \textit{Keyfile} & ~\newline \\
		\hline
	\end{tabular}
	\label{tbl:appendix:infrastruktur:cred}
	\end{table}
\end{landscape}


	%% ------------------------------------------------- Glossar ------------------------------------------------------------------
	% Glossar
	%\printglossary[style=altlist,title=Glossar]
	% Abk"urzungen
	%\deftranslation[to=German]{Acronyms}{Abk\"urzungsverzeichnis}
	%\printglossary[type=\acronymtype,style=long]
	% Symbole
	%\printglossary[type=symbolslist,style=long]
	% Englisch-Deutsch
	%\printglossary[type= translationlist, style=long]

	%% ------------------------------------------------- Literaturverzeichnis -----------------------------------------------------
	% Literaturverzeichnis
	\bibliography{bib/bib}

	%% ------------------------------------------------- Index --------------------------------------------------------------------
	% Index - Schlagworte
	\printindex
	% Index - Namensverzeichnis
	%\printindex[name]
	% Eigenst"andigkeitserkl"arung
	\newpage{}
	\pagestyle{empty}
	%\pagenumbering{} 
	\selectlanguage{ngerman}
	\section*{Erkl"arung}
	Hiermit versichere ich, dass ich diese Arbeit selbstst"andig verfasst und keine anderen als die angegebenen Quellen und Hilfsmittel benutzt habe. Au"serdem versichere ich, dass ich die allgemeinen Prinzipien wissenschaftlicher Arbeit und Ver"offentlichung, wie sie in den Leitlinien guter wissenschaftlicher Praxis der Carl von Ossietzky Universit"at Oldenburg festgelegt sind, befolgt habe.
	
	\vspace{1cm}
	Oldenburg, den \today
	
	\vspace{0.5cm}
	\rule{6cm}{0.5pt}
	
	(Muhammad Ekbal Ahmad)
\end{document}
%%%  =============================================== Document END =====================================================================
